\graphicspath{{Parti/P01-Preliminari/A04-Vettori_Polari_Assiali_Liberi_Applicati/Figure/}}
%	
\chapter{Vettori geometrici: liberi, applicati, polari e assiali}\label{app:vettori_geometrici}




\begin{sommario}
	Il capitolo presenta i fondamenti della teoria vettoriale necessari per	la meccanica strutturale. Dopo aver introdotto il concetto di vettore	libero come classe di equivalenza di segmenti orientati equipollenti,	si richiamano le operazioni fondamentali dell'algebra vettoriale ---	prodotto scalare, prodotto vettoriale, doppio prodotto vettoriale,	prodotto misto e prodotto tensoriale --- illustrandone il significato	geometrico e le proprietà algebriche. Si introduce quindi il concetto
	di vettore applicato e si definiscono il momento polare rispetto a un	polo e il momento assiale rispetto a una retta, discutendone le	proprietà di invarianza. La legge di variazione del momento al variare	del polo mostra come il momento polare sia un oggetto affine, non	libero, e mette in luce un invariante scalare che si ritrova nello studio dei sistemi di forze e dei campi di spostamento rigido.	Il capitolo si conclude con la distinzione fra vettori polari e vettori	assiali, fondata sul loro diverso comportamento sotto riflessioni del	sistema di riferimento, e con il caso piano, in cui il momento si	riduce a uno scalare con segno.
\end{sommario}

%-------------------------------------------------------------------------------------------------------------------
\section{Vettori liberi e algebra vettoriale}
%-------------------------------------------------------------------------------------------------------------------

Un \textit{vettore libero}\index{vettore!libero|textbf} è completamente determinato dal suo modulo,
dalla sua direzione e dal suo verso, indipendentemente dal punto in cui
lo si immagina applicato: formalmente, è una classe di equivalenza di
segmenti orientati equipollenti. Rispetto a una base ortonormale destra
$\{\ab_x, \ab_y, \ab_z\}$ fissata una volta per tutte, ogni vettore
geometrico $\vb \in \Uc$ è rappresentato in modo univoco dal vettore
colonna delle sue componenti:
\begin{equation}
	\vb = v_x\,\ab_x + v_y\,\ab_y + v_z\,\ab_z
	\quad\longleftrightarrow\quad
	\bm{v} = \begin{Bmatrix} v_x \\ v_y \\ v_z \end{Bmatrix} \in \mathbb{R}^3.
\end{equation}
La corrispondenza $\vb \leftrightarrow \bm{v}$ è un isomorfismo che
permette di tradurre operazioni geometriche in operazioni algebriche sulle
componenti. Nei paragrafi seguenti si richiamano le operazioni fondamentali
dell'algebra vettoriale, che saranno impiegate sistematicamente nel testo.


\subsection{Somma e prodotto per uno scalare}

La somma di due vettori $\ub, \vb \in \Uc$ si costruisce con la regola
del parallelogramma, o equivalentemente con il metodo punta-coda
(\FIGref{f:somma_vettoriale}). 
In componenti:
\begin{equation}
	\ub + \vb = (u_x+v_x)\,\ab_x + (u_y+v_y)\,\ab_y + (u_z+v_z)\,\ab_z.
\end{equation}
La somma $\ub + \vb$ si costruisce geometricamente in due modi	equivalenti (\FIGref{f:somma_vettoriale}): con la \textit{regola del 	parallelogramma}, ponendo i due vettori con la stessa origine e	prendendo la diagonale, oppure con il \textit{metodo punta-coda}, ponendo l'origine di $\vb$ sulla punta di $\ub$ e congiungendo	l'origine di $\ub$ con la punta di $\vb$. La differenza $\ub - \vb$	si riconduce alla somma con l'opposto, $\ub + (-\vb)$; equivalentemente, quando i due vettori sono applicati nella stessa origine, il vettore	differenza va dalla punta del sottraendo $\vb$ alla punta del minuendo	$\ub$.


\begin{figure}[htbp]
	\centering
	\tikzset{midarrow/.style={postaction={decorate},
			decoration={markings, mark=at position #1 with {\arrow{Stealth}}}},
		midarrow/.default=0.55}
	\begin{minipage}[b]{0.46\textwidth}\centering
		\begin{tikzpicture}[>=Stealth, line width=0.8pt, scale=1.0]
			% ===== (a) SOMMA =====
			\coordinate (O) at (0,0);
			\coordinate (A) at (3,0);   % punta di u
			\coordinate (B) at (1,2);   % punta di v
			\coordinate (C) at (4,2);   % punta di u+v
			% lati traslati (punta-coda): freccia a meta' tratto
			\draw[thick, midarrow] (B) -- (C);   % u traslato
			\draw[thick, midarrow] (A) -- (C);   % v traslato
			% vettori dalla stessa origine
			\draw[->, thick] (O) -- (A) node[midway, below=1pt] {$\ub$};
			\draw[->, thick] (O) -- (B) node[midway, above left=-3pt] {$\vb$};
			% diagonale = somma
			\draw[thick, midarrow] (O) -- (C) node[midway, sloped, above=2pt] {$\ub+\vb$};
		\end{tikzpicture}\\[0.8ex]
		\emph{(a) somma}
	\end{minipage}\hfill
	\begin{minipage}[b]{0.46\textwidth}\centering
		\begin{tikzpicture}[>=Stealth, line width=0.8pt, scale=1.0]
			% ===== (b) DIFFERENZA =====
			\coordinate (O) at (0,0);
			\coordinate (A) at (3,0);   % punta di u
			\coordinate (B) at (1,2);   % punta di v
			\coordinate (C) at (4,2);
			% parallelogramma di contesto (lati di completamento)
			\draw[gray!55, thin] (A) -- (C);
			\draw[gray!55, thin] (B) -- (C);
			% vettori u, v
			\draw[->, thick] (O) -- (A) node[midway, below=1pt] {$\ub$};
			\draw[->, thick] (O) -- (B) node[midway, above left=-3pt] {$\vb$};
			% differenza: dalla punta di v alla punta di u
			\draw[thick, midarrow] (B) -- (A) node[midway, sloped, above=2pt] {$\ub-\vb$};
		\end{tikzpicture}\\[0.8ex]
		\emph{(b) differenza}
	\end{minipage}
	\caption{Somma e differenza di due vettori.}
	\label{f:somma_vettoriale}
\end{figure}


Il prodotto di un vettore $\vb$ per uno scalare $\alpha \in \mathbb{R}$
ne scala il modulo di un fattore $|\alpha|$ e ne inverte il verso se
$\alpha < 0$ (\FIGref{f:prodotto_per_uno_scalare}):
\begin{equation}
	\alpha\vb = \alpha v_x\,\ab_x + \alpha v_y\,\ab_y + \alpha v_z\,\ab_z.
\end{equation}

\begin{figure}[htbp]
	\centering
	
	\begin{tikzpicture}[>=Stealth, line width=0.8pt, scale=0.80]
		% ===================== sinistra: alpha = 2 =====================
		\coordinate (OL) at (0,0);
		% v
		\draw[->, thick] (OL) -- ++(0.85,1.6) node[midway, above left=-3pt] {$\vb$};
		% alpha v = 2v (parallelo, traslato a destra, lunghezza doppia)
		\draw[->, thick] ($(OL)+(0.8,-0.2)$) -- ++(1.7,3.2)
		node[midway, right=2pt] {$\alpha\vb = 2\vb$};
		%\node at ($(OL)+(1.15,3.55)$) {$\alpha = 2$};
		
		% ===================== destra: alpha = -1 =====================
		\coordinate (OR) at (5.2,0);
		% v
		\draw[->, thick] (OR) -- ++(0.85,1.6) node[midway, above left=-3pt] {$\vb$};
		% alpha v = -v (parallelo, traslato a destra, verso opposto)
		\draw[->, thick] ($(OR)+(2.0,1.85)$) -- ++(-0.85,-1.6)
		node[midway, right=2pt] {$\alpha\vb = -\vb$};
		%\node at ($(OR)+(1.0,3.55)$) {$\alpha = -1$};
	\end{tikzpicture}
	\caption{Prodotto di un vettore $\vb$ per uno scalare:
		$\alpha = 2$ (sinistra) e $\alpha = -1$ (destra).}
	\label{f:prodotto_per_uno_scalare}
\end{figure}




\subsection{Prodotto scalare}\index{prodotto scalare|textbf}\nomenclature[O]{$\ub\cdot\vb$}{prodotto scalare}

Il prodotto scalare di $\ub, \vb \in \Uc$ è lo scalare
\begin{equation}\label{eq:prod_scal}
	\ub \cdot \vb \coloneqq |\ub|\,|\vb|\cos\alpha,
\end{equation}
dove $\alpha \in [0,\pi]$ è l'angolo convesso fra i due vettori
(\FIGref{f:prodotto_scalare}). In componenti, rispetto alla base
ortonormale:
\begin{equation}
	\ub \cdot \vb = u_x v_x + u_y v_y + u_z v_z = \bm{u}^\top\bm{v}.
\end{equation}
Proprietà fondamentali:
\begin{enumerate}[label=(\roman*)]
	\item $\ub \cdot \vb = 0$ se e solo se $\ub \perp \vb$;
	\item $\vb \cdot \vb = |\vb|^2$, da cui $|\vb| = \sqrt{\bm{v}^\top\bm{v}}$;
	\item la proiezione orientata di $\vb$ su $\ub$ è lo scalare
	$|\vb|\cos\alpha = (\ub \cdot \vb)/|\ub|$, sicché
	$\ub \cdot \vb = |\ub|\,\mathrm{proj}_{\ub}\vb
	= |\vb|\,\mathrm{proj}_{\vb}\ub$.
\end{enumerate}
\begin{figure}[htbp]
	\centering
	% marcatore d'angolo retto in #2 fra i raggi verso #1 e #3 (orientamento qualsiasi)
	\providecommand{\rightangle}[3]{%
		\coordinate (ra1) at ($(#2)!2.4mm!(#1)$);
		\coordinate (ra3) at ($(#2)!2.4mm!(#3)$);
		\draw[gray!55, line width=0.5pt] (ra1) -- ($(ra1)+(ra3)-(#2)$) -- (ra3);}
	\begin{minipage}[b]{0.48\textwidth}\centering
		\begin{tikzpicture}[>=Stealth, line width=0.8pt, scale=1.0]
			% ===== (a) proiezione di v su u =====
			\coordinate (O) at (0,0);
			\coordinate (U) at (3.2,0);
			\coordinate (V) at (40:2.8);
			\coordinate (F) at ($(O)!(V)!(U)$);   % piede sulla retta di u
			\draw[gray!55, thin] (V) -- (F);
			\rightangle{O}{F}{V}
			\draw[->, thick] (O) -- (U) node[right=2pt] {$\ub$};
			\draw[->, thick] (O) -- (V) node[above=2pt, midway] {$\vb$};
			\pic[draw, thin, angle radius=9mm, angle eccentricity=1.28,
			"$\alpha$"{font=\small}] {angle = U--O--V};
			% quote orizzontali sotto l'asse di u
			\draw[gray!55, thin] (O) -- (0,-1.5);
			\draw[gray!55, thin] (F) -- ($(F)+(0,-0.75)$);
			\draw[gray!55, thin] (U) -- ($(U)+(0,-1.5)$);
			\draw[<->, thin] (0,-0.6) -- ($(F)+(0,-0.6)$)
			node[midway, above, font=\small] {$|\vb|\cos\alpha$};
			\draw[<->, thin] (0,-1.35) -- ($(U)+(0,-1.35)$)
			node[midway, above, font=\small] {$|\ub|$};
			\fill (O) circle (1.2pt);
		\end{tikzpicture}\\[0.6ex]
		\emph{(a) proiezione di $\vb$ su $\ub$}
	\end{minipage}\hfill
	\begin{minipage}[b]{0.48\textwidth}\centering
		\begin{tikzpicture}[>=Stealth, line width=0.8pt, scale=1.0]
			% ===== (b) proiezione di u (orizzontale) su v (inclinato) =====
			\coordinate (O) at (0,0);
			\coordinate (U) at (3.2,0);
			\coordinate (V) at (40:2.8);
			\coordinate (H) at ($(O)!(U)!(V)$);   % piede sulla retta di v
			\coordinate (Vrot) at ($(O)!1!90:(V)$); % v ruotato +90 (per le quote inclinate)
			\draw[gray!55, thin] (U) -- (H);
			\rightangle{O}{H}{U}
			\draw[->, thick] (O) -- (U) node[right=2pt] {$\ub$};
			\draw[->, thick] (O) -- (V) node[above=2pt, pos=0.5] {$\vb$};
			\pic[draw, thin, angle radius=9mm, angle eccentricity=1.28,
			"$\alpha$"{font=\small}] {angle = U--O--V};
			% quote inclinate, parallele a v (lato superiore-sinistro)
			\draw[gray!55, thin] (O) -- ($(O)+0.47*(Vrot)$);
			\draw[gray!55, thin] (H) -- ($(H)+0.23*(Vrot)$);
			\draw[gray!55, thin] (V) -- ($(V)+0.47*(Vrot)$);
			\draw[<->, thin] ($(O)+0.18*(Vrot)$) -- ($(H)+0.18*(Vrot)$)
			node[midway, sloped, above, font=\small] {$|\ub|\cos\alpha$};
			\draw[<->, thin] ($(O)+0.43*(Vrot)$) -- ($(V)+0.43*(Vrot)$)
			node[midway, sloped, above=1pt, font=\small] {$|\vb|$};
			\fill (O) circle (1.2pt);
		\end{tikzpicture}\\[0.6ex]
		\emph{(b) proiezione di $\ub$ su $\vb$}
	\end{minipage}
	\caption{Rappresentazione geometrica del prodotto scalare.}
	\label{f:prodotto_scalare}
\end{figure}


\subsection{Prodotto vettoriale}\index{prodotto vettoriale|textbf}\nomenclature[O]{$\ub\times\vb$}{prodotto vettoriale}

Il prodotto vettoriale $\ub \times \vb$ è il vettore ortogonale al piano
di $\ub$ e $\vb$, con verso dato dalla regola della mano destra e modulo
\begin{equation}\label{eq:prod_vett_modulo}
	|\ub \times \vb| \coloneqq |\ub|\,|\vb|\sin\alpha,
\end{equation}
pari all'area del parallelogramma costruito sui due vettori
(\FIGref{f:prodotto_vettoriale}). In componenti:
\begin{equation}\label{eq:prod_vett_comp}
	\ub \times \vb =
	\begin{vmatrix}
		\ab_x & \ab_y & \ab_z \\
		u_x & u_y & u_z \\
		v_x & v_y & v_z
	\end{vmatrix}
	= (u_y v_z - u_z v_y)\,\ab_x
	+ (u_z v_x - u_x v_z)\,\ab_y
	+ (u_x v_y - u_y v_x)\,\ab_z.
\end{equation}
Posto $\bm{\upomega}\coloneqq \ub \times \vb$, le tre componenti si sintetizzano	nella forma compatta
\begin{equation}\label{eq:prod_vett_ricci}
	\omega_i = \varepsilon_{ijk}\,u_j\,v_k,
\end{equation}
dove si sottintende la somma sugli indici ripetuti (convenzione di
Einstein) e $\varepsilon_{ijk}$ è il \textit{simbolo di Ricci}
(o di Levi-Civita): vale $+1$ se $(i,j,k)$ è una permutazione pari
di $(x,y,z)$, $-1$ se dispari, $0$ se due indici coincidono.


Proprietà fondamentali:
\begin{enumerate}[label=(\roman*)]
	\item antisimmetria: $\ub \times \vb = -(\vb \times \ub)$;
	\item $\ub \times \vb = \zerovb$ se e solo se $\ub$ e $\vb$
	sono paralleli.
\end{enumerate}


\begin{figure}[htbp]
	\centering
	% ===== helper condivisi da TUTTI i pannelli (definiti una volta sola) =====
	\def\ax{0.5}\def\ay{0.30}
	\providecommand{\xyz}[3]{{#2 + (#1)*\ax},{#3 + (#1)*\ay}}
	\providecommand{\rightangle}[3]{%
		\coordinate (ra1) at ($(#2)!2.2mm!(#1)$);
		\coordinate (ra3) at ($(#2)!2.2mm!(#3)$);
		\draw[gray!55, line width=0.5pt] (ra1) -- ($(ra1)+(ra3)-(#2)$) -- (ra3);}
	
	% =========================== RIGA 1 ===========================
	\begin{minipage}[b]{0.48\linewidth}
		\centering
		% --------- PANNELLO (a): vista 3D ---------
		\begin{tikzpicture}[>=Stealth, line width=0.8pt]
			\coordinate (O)  at (\xyz{0}{0}{0});
			\coordinate (U)  at (\xyz{-0.8}{2.6}{0});
			\coordinate (V)  at (\xyz{2.6}{0.7}{0});
			\coordinate (UV) at (\xyz{1.8}{3.3}{0});
			\coordinate (X)  at (\xyz{0}{0}{2.2});
			\fill[gray!7] (\xyz{-1.2}{-1.0}{0}) -- (\xyz{2.6}{-0.5}{0})
			-- (\xyz{2.6}{3.6}{0})  -- (\xyz{-1.2}{3.1}{0}) -- cycle;
			\draw[gray!40, thin] (\xyz{-1.2}{-1.0}{0}) -- (\xyz{2.6}{-0.5}{0})
			-- (\xyz{2.6}{3.6}{0})  -- (\xyz{-1.2}{3.1}{0}) -- cycle;
			\fill[gray!22] (O) -- (U) -- (UV) -- (V) -- cycle;
			\draw[gray!60, thin] (U) -- (UV) -- (V);
			\node[gray!75, font=\footnotesize] at (\xyz{0.8}{1.9}{0}) {$\mathrm{Span}\{\ub,\vb\}$};
			\draw[->, thick] (O) -- (X) node[above=1pt] {$\ub\times\vb$};
			\rightangle{X}{O}{U}
			\rightangle{X}{O}{V}
			\draw[->, thick] (O) -- (U) node[midway, below right=-2pt] {$\ub$};
			\draw[->, thick] (O) -- (V) node[midway, above=1pt] {$\vb$};
			\pic[draw, ->, thin, angle radius=6mm, angle eccentricity=1.4,	"$\alpha$"{font=\footnotesize}] {angle = U--O--V};
			\fill (O) circle (1.2pt);
		\end{tikzpicture}\\[4pt]
		{\small (a)}
	\end{minipage}\hfill
	\begin{minipage}[b]{0.48\linewidth}
		\centering
		% --------- PANNELLO (b): area nel piano ---------
		\begin{tikzpicture}[>=Stealth, line width=0.8pt]
			\coordinate (Od)  at (0,0);
			\coordinate (Ud)  at (2.6,0);
			\coordinate (Vd)  at (40:1.9);
			\coordinate (UVd) at ($(Ud)+(Vd)$);
			\coordinate (Fd)  at ($(Od)!(Vd)!(Ud)$);
			\fill[gray!14] (Od) -- (Ud) -- (UVd) -- (Vd) -- cycle;
			\draw[gray!60, thin] (Ud) -- (UVd) -- (Vd);
			\draw[gray!55, thin, dashed] (Vd) -- (Fd);
			\rightangle{Od}{Fd}{Vd}
			
			\node[font=\footnotesize, right=1pt] at ($(Vd)!0.22!(Fd)$) {$|\vb|\sin\alpha$};
			\draw[->, thick] (Od) -- (Ud) node[midway, below=2pt]{$\ub$};
			\draw[->, thick] (Od) -- (Vd) node[pos=0.62, above left=-3pt]{$\vb$};
			\pic[draw, ->, thin, angle radius=6mm, angle eccentricity=1.4, "$\alpha$"{font=\footnotesize}] {angle = Ud--Od--Vd};
			\node[font=\footnotesize] at (3.0,0.62) {$|\ub\times\vb|$};
			\draw[gray!55, thin] (Od) -- ($(Od)+(0,-0.95)$);
			\draw[gray!55, thin] (Ud) -- ($(Ud)+(0,-0.95)$);
			\draw[<->, thin] ($(Od)+(0,-0.8)$) -- ($(Ud)+(0,-0.8)$)
			node[midway, below=1pt, font=\footnotesize]{$|\ub|$};
		\end{tikzpicture}\\[4pt]
		{\small (b)}
	\end{minipage}
	
	\vspace{1.6em}
	
	% =========================== RIGA 2 ===========================
	\begin{minipage}[b]{0.48\linewidth}
		\centering
		% --------- PANNELLO (b) sinistra: u x v uscente ---------
		\begin{tikzpicture}[>=Stealth, line width=0.8pt]
			\coordinate (Ob) at (0,0);
			\coordinate (Ub) at (1.7,0);
			\coordinate (Vb) at (42:1.55);
			\draw (Ob) circle (0.17); \fill (Ob) circle (1.3pt);
			\draw[->, thick] ($(Ob)+(0.17,0)$) -- (Ub) node[right=1pt]{$\ub$};
			\draw[->, thick] ($(Ob)+(42:0.17)$) -- (Vb) node[above right=-2pt]{$\vb$};
			\pic[draw, ->, thin, angle radius=6mm, angle eccentricity=1.35,	"$\alpha$"{font=\footnotesize}] {angle = Ub--Ob--Vb};
			\node[left=4pt, font=\footnotesize] at (Ob) {$\ub\times\vb$};
		\end{tikzpicture}\\[4pt]
		{\small (c)\quad $\ub\times\vb$ uscente ($\odot$)}
	\end{minipage}\hfill
	\begin{minipage}[b]{0.48\linewidth}
		\centering
		% --------- PANNELLO (b) destra: v x u entrante ---------
		\begin{tikzpicture}[>=Stealth, line width=0.8pt]
			\coordinate (Oc) at (0,0);
			\coordinate (Uc) at (1.7,0);
			\coordinate (Vc) at (42:1.55);
			\draw (Oc) circle (0.17);
			\draw (-0.12,-0.12) -- (0.12,0.12);
			\draw (-0.12,0.12) -- (0.12,-0.12);
			\draw[->, thick] ($(Oc)+(0.17,0)$) -- (Uc) node[right=1pt]{$\ub$};
			\draw[->, thick] ($(Oc)+(42:0.17)$) -- (Vc) node[above right=-2pt]{$\vb$};
			\pic[draw, <-, thin, angle radius=6mm, angle eccentricity=1.35,
			"$\alpha$"{font=\footnotesize}] {angle = Uc--Oc--Vc};
			\node[left=4pt, font=\footnotesize] at (Oc) {$\vb\times\ub$};
		\end{tikzpicture}\\[4pt]
		{\small (d)\quad $\vb\times\ub$ entrante ($\otimes$)}
	\end{minipage}
	
	\caption{Rappresentazione geometrica del prodotto vettoriale $\ub\times\vb$.
		In (a) il vettore $\ub\times\vb$ è ortogonale al piano
		$\mathrm{Span}\{\ub,\vb\}$ e orientato secondo la regola della mano
		destra; il suo modulo eguaglia l'area del parallelogramma costruito
		su $\ub$ e $\vb$. In (b) tale area è letta come prodotto della base
		$|\ub|$ per l'altezza $|\vb|\sin\alpha$, da cui
		$|\ub\times\vb| = |\ub|\,|\vb|\sin\alpha$. In (c) e (d) è illustrata
		l'antisimmetria $\ub\times\vb = -(\vb\times\ub)$: a parità di $\ub$,
		$\vb$ e di angolo $\alpha$, il prodotto cambia verso --- uscente dal
		piano ($\odot$) per $\ub\times\vb$, entrante ($\otimes$) per
		$\vb\times\ub$ --- coerentemente col senso di rotazione che porta il
		primo fattore sul secondo.}
	\label{f:prodotto_vettoriale}
\end{figure}


\begin{oss}[\textbf{Forma matriciale}]
	Il prodotto vettoriale può essere scritto come prodotto
	matrice-vettore tramite la matrice antisimmetrica associata a $\ub$:
	\begin{equation}\label{eq:prod_vett_matrice}
		\bm{u} \times \bm{v}
		= \begin{bmatrix}
			0   & -u_z &  u_y \\
			u_z &  0   & -u_x \\
			-u_y &  u_x &  0
		\end{bmatrix}
		\begin{Bmatrix} v_x \\ v_y \\ v_z \end{Bmatrix}.
	\end{equation}
	Questa rappresentazione sarà impiegata nel calcolo del momento e nella
	formulazione matriciale della cinematica infinitesima.
\end{oss}


\subsection{Doppio prodotto vettoriale}

Il doppio prodotto vettoriale si esprime mediante l'identità di Grassmann:
\begin{equation}\label{eq:doppio_vett}
	(\ub \times \vb) \times \wb	= (\ub \cdot \wb)\,\vb - (\vb \cdot \wb)\,\ub,
\end{equation}
che mostra come il risultato giaccia sempre nel piano di $\ub$ e $\vb$
(\FIGref{f:doppio_prodotto_vettoriale}). Il prodotto vettoriale non è
associativo: in generale
$(\ub \times \vb) \times \wb \neq \ub \times (\vb \times \wb)$,
poiché la seconda forma vale
$\ub \times (\vb \times \wb) = (\ub \cdot \wb)\,\vb - (\ub \cdot \vb)\,\wb$,
che giace nel piano di $\vb$ e $\wb$.

\begin{figure}[!htbp]
	\centering
	\def\ax{0.5}\def\ay{0.30}
	\providecommand{\xyz}[3]{{#2 + (#1)*\ax},{#3 + (#1)*\ay}}
	\providecommand{\rightangle}[3]{%
		\coordinate (ra1) at ($(#2)!2.2mm!(#1)$);
		\coordinate (ra3) at ($(#2)!2.2mm!(#3)$);
		\draw[gray!55, line width=0.5pt] (ra1) -- ($(ra1)+(ra3)-(#2)$) -- (ra3);}
	\begin{tikzpicture}[>=Stealth, line width=0.8pt, scale=0.80]
		% u, v nel piano base (z=0); u x v verticale; (u x v) x w nel piano base
		\coordinate (O)  at (\xyz{0}{0}{0});
		\coordinate (A)  at (\xyz{-0.8}{3.0}{0});   % u
		\coordinate (B)  at (\xyz{2.8}{0.9}{0});    % v
		\coordinate (UV) at (\xyz{2.0}{3.9}{0});    % u+v
		\coordinate (X)  at (\xyz{0}{0}{2.6});      % u x v (verticale)
		\coordinate (W)  at (\xyz{1.0}{1.2}{1.8});  % w
		\coordinate (P)  at (\xyz{3.8}{-3.2}{0});   % (u x v) x w  (combinazione di u,v)
		\coordinate (AP) at (\xyz{1.0}{1.2}{4.4});  % apice = (u x v)+w
		
		% piano base (su u,v), esteso per contenere anche (u x v) x w
		\coordinate (Q1) at (\xyz{-1.2}{-5.7}{0});   % era -4.3
		\coordinate (Q2) at (\xyz{4.3}{-4.9}{0});    % era -3.5
		\coordinate (Q3) at (\xyz{4.3}{5.5}{0});
		\coordinate (Q4) at (\xyz{-1.2}{4.2}{0});
		\fill[gray!6] (Q1) -- (Q2) -- (Q3) -- (Q4) -- cycle;
		\draw[gray!40, thin] (Q1) -- (Q2) -- (Q3) -- (Q4) -- cycle;
		% parallelogramma costruito su u e v
		\fill[gray!18] (O) -- (A) -- (UV) -- (B) -- cycle;
		\draw[gray!60, thin] (A) -- (UV) -- (B);
		% piano di (u x v, w): lati ausiliari verso l'apice
		\draw[gray!55, thin] (X) -- (AP) -- (W);
		
		% angolo retto: (u x v) x w perpendicolare a w
		\rightangle{W}{O}{P}
		\rightangle{X}{O}{P}
		
		% vettori base
		\draw[->, thick] (O) -- (A) node[midway, above, xshift=14pt] {$\ub$};
		\draw[->, thick] (O) -- (B) node[midway, above left=-2pt, xshift=6pt] {$\vb$};
		% u x v, w, (u x v) x w
		\draw[->, thick] (O) -- (X) node[above left=-1pt] {$\ub\times\vb$};
		\draw[->, thick] (O) -- (W) node[midway, above=4pt] {$\wb$};
		\draw[->, thick] (O) -- (P) node[below left=-6pt,  yshift=-8pt] {$(\ub\times\vb)\times\wb$};
		
		\fill (O) circle (1.2pt);
	\end{tikzpicture}
	\caption{Il doppio prodotto vettoriale $(\ub\times\vb)\times\wb$
		giace nel piano di $\ub$ e $\vb$.}
	\label{f:doppio_prodotto_vettoriale}
\end{figure}


\subsection{Prodotto misto}\index{prodotto misto|textbf}

Il prodotto misto di tre vettori $\ub, \vb, \wb \in \Uc$ è lo scalare
\begin{equation}\label{eq:prod_misto}
	(\ub \times \vb) \cdot \wb
	= \begin{vmatrix}
		u_x & u_y & u_z \\
		v_x & v_y & v_z \\
		w_x & w_y & w_z
	\end{vmatrix}.
\end{equation}
Il suo valore assoluto è il volume del parallelepipedo costruito sui tre
vettori (\FIGref{f:prodotto_misto}); il segno è positivo se $\ub, \vb,
\wb$ formano una terna destrorsa, negativo altrimenti. Il prodotto misto
è nullo se e solo se i tre vettori sono complanari, ed è invariante per
permutazioni cicliche:
\begin{equation}
	(\ub \times \vb) \cdot \wb
	= (\vb \times \wb) \cdot \ub
	= (\wb \times \ub) \cdot \vb,
\end{equation}
mentre cambia segno per scambio di due fattori qualsiasi.


\begin{figure}[!htbp]
	\centering
	% --- helper locali (blocco autonomo) ---
	\def\ax{0.5}\def\ay{0.30}
	\providecommand{\xyz}[3]{{#2 + (#1)*\ax},{#3 + (#1)*\ay}}
	\providecommand{\rightangle}[3]{%
		\coordinate (ra1) at ($(#2)!2.4mm!(#1)$);
		\coordinate (ra3) at ($(#2)!2.4mm!(#3)$);
		\draw[gray!55, line width=0.5pt] (ra1) -- ($(ra1)+(ra3)-(#2)$) -- (ra3);}
	\begin{tikzpicture}[>=Stealth, line width=0.8pt, scale=1.0]
		% ---- vertici del parallelepipedo (x=profondita', y=oriz, z=vert) ----
		\coordinate (O)   at (\xyz{0}{0}{0});
		\coordinate (A)   at (\xyz{-0.5}{3.2}{0});      % O+u
		\coordinate (B)   at (\xyz{2.8}{0.6}{0});       % O+v
		\coordinate (C)   at (\xyz{0.5}{0.7}{2.4});     % O+w
		\coordinate (AB)  at (\xyz{2.3}{3.8}{0});       % u+v
		\coordinate (AC)  at (\xyz{0.0}{3.9}{2.4});     % u+w
		\coordinate (BC)  at (\xyz{3.3}{1.3}{2.4});     % v+w
		\coordinate (ABC) at (\xyz{2.8}{4.5}{2.4});     % u+v+w
		% asse di u x v (verticale) e piede H
		\coordinate (X)   at (\xyz{0}{0}{2.95});        % punta di u x v
		\coordinate (H)   at (\xyz{0}{0}{2.4});         % piede (quota di C)
		
		% ---- spigoli del parallelepipedo (sottili) ----
		\begin{scope}[gray!60, line width=0.6pt]
			% base
			\draw (A) -- (AB) -- (B);
			% faccia superiore
			\draw (C) -- (AC) -- (ABC) -- (BC) -- cycle;
			% spigoli laterali (lungo w)
			\draw (A) -- (AC);
			\draw (AB) -- (ABC);
			\draw (B) -- (BC);
		\end{scope}
		
		% ---- altezza: u x v, piede H, proiezione orizzontale H--C ----
		\draw[gray!55, thin, dashed] (H) -- (C);
		\draw[->, thick] (O) -- (X) node[above=1pt] {$\ub\times\vb$};
		\node[left=1pt] at (H) {$H$};
		\rightangle{O}{H}{C}
		
		% ---- i tre vettori u, v, w (spessi) ----
		\draw[->, thick] (O) -- (A) node[midway, below=1pt, xshift=6pt] {$\ub$};
		\draw[->, thick] (O) -- (B) node[midway, below, xshift=12pt, yshift=4pt] {$\vb$};
		\draw[->, thick] (O) -- (C) node[midway, right=1pt] {$\wb$};
		
		% ---- punti ----
		\fill (O) circle (1.4pt) node[left=2pt] {$O$};
		\fill (A) circle (1.4pt) node[below right=-2pt] {$A$};
		\fill (B) circle (1.4pt) node[above left=-1pt] {$B$};
		\fill (C) circle (1.4pt) node[above left=-2pt] {$C$};
	\end{tikzpicture}
	\caption{Il valore assoluto del prodotto misto è il volume del
		parallelepipedo costruito su $\ub$, $\vb$ e $\wb$.}
	\label{f:prodotto_misto}
\end{figure}


\subsection{Prodotto tensoriale}\index{prodotto tensoriale|textbf}\nomenclature[O]{$\ub\otimes\vb$}{prodotto tensoriale}

Il prodotto tensoriale di due vettori $\ub, \vb \in \Uc$, indicato con
$\ub \otimes \vb$, è il tensore del secondo ordine definito dalla
condizione
\begin{equation}
	(\ub \otimes \vb)\,\wb \coloneqq (\vb \cdot \wb)\,\ub
	\qquad \forall\,\wb \in \Uc.
\end{equation}
In componenti, la matrice associata ha elementi
$(\ub \otimes \vb)_{ij} = u_i v_j$, ovvero $\bm{A} = \bm{u}\bm{v}^\top$.
Un caso di particolare rilievo è il prodotto tensoriale di un versore
$\ab$ per sé stesso: il tensore $\ab \otimes \ab$ è il
\textit{proiettore} ortogonale sulla direzione di $\ab$, poiché
\begin{equation}
	(\ab \otimes \ab)\,\vb = (\ab \cdot \vb)\,\ab = \mathrm{proj}_{\ab}\vb\cdot\ab,
\end{equation}
estrae cioè la componente di $\vb$ lungo $\ab$. Il tensore complementare
$\Ib - \ab \otimes \ab$ proietta invece sul piano ortogonale ad $\ab$.


%-------------------------------------------------------------------------------------------------------------------
\section{Vettori applicati}\label{sec:vettori_applicati}
%-------------------------------------------------------------------------------------------------------------------

Un \textit{vettore applicato}\index{vettore!applicato|textbf} è una coppia ordinata $(P,\,\vb)$, dove
$P \in \Ec$ è il \textit{punto di applicazione} e $\vb \in \Uc$ è un
vettore libero. Mentre un vettore libero è individuato da tre parametri
(le componenti), un vettore applicato ne richiede sei: tre per la
posizione di $P$ e tre per $\vb$. La retta passante per $P$ e parallela
a $\vb$ si chiama \textit{linea d'azione}. La distinzione fra vettori
liberi e applicati non è formale: forze, spostamenti e velocità sono
vettori applicati, perché il loro effetto fisico dipende dal punto in
cui agiscono.

\subsection{Momento polare}\label{sec:momento_polare}

Dato un vettore applicato $(P,\,\vb)$ e un punto $O \in \Ec$ detto
\textit{polo}\index{polo}, si definisce \textit{momento polare}\index{momento!polare|textbf}\nomenclature[G]{$\bm{\upmu}$}{momento polare di un vettore applicato} di $\vb$ rispetto
a $O$ il vettore
\begin{equation}\label{eq:momento_polare}
	\bm{\upmu}(O) \coloneqq \overrightarrow{OP} \times \vb.
\end{equation}
Il vettore $\bm{\upmu}(O)$ è perpendicolare al piano di
$\overrightarrow{OP}$ e $\vb$, con verso dato dalla regola della mano
destra. Il suo modulo è
\begin{equation}
	|\bm{\upmu}(O)| = |\overrightarrow{OP}|\,|\vb|\sin\alpha = |\vb|\,b,
\end{equation}
dove $\alpha$ è l'angolo fra $\overrightarrow{OP}$ e $\vb$, e $b$ è la
distanza della linea d'azione dal polo, detto \textit{braccio}\index{braccio|textbf}
(\FIGref{f:momento_polare}).

\begin{figure}[!htbp]
	\centering
	\def\ax{0.5}\def\ay{0.30}
	\providecommand{\xyz}[3]{{#2 + (#1)*\ax},{#3 + (#1)*\ay}}
	\providecommand{\rightangle}[3]{%
		\coordinate (ra1) at ($(#2)!2.4mm!(#1)$);
		\coordinate (ra3) at ($(#2)!2.4mm!(#3)$);
		\draw[line width=0.5pt] (ra1) -- ($(ra1)+(ra3)-(#2)$) -- (ra3);}
	\tikzset{midarrow/.style={postaction={decorate},decoration={markings, mark=at position 0.56 with {\arrow{Stealth}}}}}
	%
	\begin{minipage}[b]{0.46\textwidth}\centering
		\begin{tikzpicture}[>=Stealth, line width=0.8pt, scale=1.0]
			% ===== PANNELLO (a) =====
			\draw[gray!60] (\xyz{-1.7}{-1.2}{0}) -- (\xyz{-1.7}{3.50}{0})-- (\xyz{3.0}{3.50}{0}) -- (\xyz{3.0}{-1.2}{0}) -- cycle;
			\node[gray!75] at (\xyz{-1.2}{-1.0}{0}) {$\pi$};
			\coordinate (O)  at (\xyz{0}{0}{0});
			\coordinate (Pk) at (\xyz{0.5}{2.0}{0});
			\coordinate (Ff) at (\xyz{1.7}{2.3}{0});
			\coordinate (Of) at (\xyz{1.2}{0.3}{0});
			\coordinate (MuB) at (\xyz{0}{0}{1.25});
			\coordinate (MuC) at (\xyz{0}{0}{1.47});
			\draw[gray!55, thin] (O) -- (Of) -- (Ff) -- (Pk) -- cycle;
			%\draw[->, thin] (O) -- (Pk);
			%\node[font=\small, below] at ($(O)!0.5!(Pk)$) {$\overrightarrow{OP}$};
			
			\draw[->,thick] (O) -- (Pk) node[font=\small, midway, sloped, below] {$\overrightarrow{OP}$};
			
			\draw[->, thick] (Pk) -- (Ff) node[right=-2pt] {$\vb$};
			\draw[->>, thick] (O) -- (MuB) node[left=1pt, yshift=-10pt] {$\bm{\upmu}(O)$};
			\begin{scope}[gray!55]
				\rightangle{MuB}{O}{Of}
				\rightangle{MuB}{O}{Pk}
			\end{scope}
			\fill (O)  circle (1.5pt) node[left=2pt] {$O$};
			\fill (Pk) circle (1.5pt) node[below right=-1pt] {$P$};
		\end{tikzpicture}\\[0.8ex]
		(a)
	\end{minipage}\hfill
	\begin{minipage}[b]{0.50\textwidth}\centering
		\begin{tikzpicture}[>=Stealth, line width=0.8pt, scale=1.0]
			% ===== PANNELLO (b) =====
			\draw[gray!60] (\xyz{-1.7}{-1.2}{0}) -- (\xyz{-1.7}{3.50}{0})-- (\xyz{3.0}{3.50}{0}) -- (\xyz{3.0}{-1.2}{0}) -- cycle;
			\node[gray!75] at (\xyz{-1.2}{-1.0}{0}) {$\pi$};
			\coordinate (O)  at (\xyz{0}{0}{0});
			\coordinate (Pk) at (\xyz{0.5}{2.0}{0});
			\coordinate (Ff) at (\xyz{1.7}{2.3}{0});
			\coordinate (Of) at (\xyz{1.2}{0.3}{0});
			\coordinate (MuB) at (\xyz{0}{0}{1.25});
			% piede della perpendicolare da O alla retta d'azione (calcolato nel piano)
			\coordinate (Foot) at (\xyz{-0.4412}{1.7647}{0});
			% retta d'azione di v (estesa)
			\coordinate (LAa) at (\xyz{-1.05}{1.6150}{0});
			\coordinate (LAb) at (\xyz{2.30}{2.4500}{0});
			% marcatore angolo retto al piede (vertici calcolati nel piano)
			\coordinate (raA) at (\xyz{-0.3878}{1.5513}{0});
			\coordinate (raB) at (\xyz{-0.1744}{1.6046}{0});
			\coordinate (raC) at (\xyz{-0.2277}{1.8181}{0});
			
			\begin{scope}[gray!55]
				\rightangle{MuB}{O}{Of}
				\rightangle{MuB}{O}{Pk}
			\end{scope}
			
			% --- retta d'azione di v (dietro a tutto) ---
			\draw[gray!55, thin] (LAa) -- (LAb);
			
			% --- parallelogramma riempito (riuso da (a)) ---
			\filldraw[gray!12, draw=gray!55, thin] (O) -- (Of) -- (Ff) -- (Pk) -- cycle;
			
			% --- braccio b + angolo retto al piede ---
			\draw[thin] (O) -- (Foot);
			\node[font=\small, below=0pt] at ($(O)!0.5!(Foot)$) {$b$};
			\draw[line width=0.5pt] (raA) -- (raB) -- (raC);
			
			% --- OP ---
			%				\draw[->, thin] (O) -- (Pk);
			%				\node[font=\small, above left=-3pt] at ($(O)!0.5!(Pk)$) {$\overrightarrow{OP}$};
			
			\draw[->,thin] (O) -- (Pk) node[font=\small, midway, sloped, above] {$\overrightarrow{OP}$};
			
			% --- v ---
			\draw[->, thick] (Pk) -- (Ff) node[above=-2pt] {$\vb$};
			
			% --- angolo alpha fra OP e v in P ---
			\coordinate (Pext) at ($(O)!1.33!(Pk)$);
			\pic[draw, thin, angle radius=6.5mm, angle eccentricity=1.5, "$\alpha$"{font=\small}] {angle = Pext--Pk--Ff};
			
			\coordinate (Pext) at ($(O)!1.45!(Pk)$);   % punto oltre P lungo OP
			\draw[gray!55, thin] (Pk) -- (Pext);        % prolungamento grigio
			
			% --- momento mu(O) ---
			\draw[->>, thick] (O) -- (MuB) node[left=1pt, yshift=-10pt] {$\bm{\upmu}(O)$};
			
			% --- punti ---
			\fill (O)  circle (1.5pt) node[left=2pt] {$O$};
			\fill (Pk) circle (1.5pt) node[below right=-1pt] {$P$};
		\end{tikzpicture}\\[0.8ex]
		(b)
	\end{minipage}
	\caption{Momento polare di un vettore applicato: (a)~definizione
		come prodotto vettoriale $\protect\overrightarrow{OP}\times\vb$;
		(b)~calcolo come prodotto del modulo $|\vb|$ per il
		braccio~$b$.}
	\label{f:momento_polare}
\end{figure}


In componenti, rispetto a un sistema cartesiano con origine in $O$:
\begin{equation}\label{eq:momento_matrice}
	\begin{Bmatrix} \mu_x(O) \\ \mu_y(O) \\ \mu_z(O) \end{Bmatrix}
	=
	\begin{bmatrix}
		0    & -z(P) &  y(P) \\
		z(P) &  0    & -x(P) \\
		-y(P) &  x(P) &  0
	\end{bmatrix}
	\begin{Bmatrix} v_x \\ v_y \\ v_z \end{Bmatrix},
\end{equation}
dove $x(P)$, $y(P)$, $z(P)$ sono le coordinate di $P$ rispetto a $O$ (\FIGref{f:componenti_momento_polare}).


\begin{figure}[!htbp]
	\centering
	\def\ax{0.5}\def\ay{0.30}
	\providecommand{\xyz}[3]{{#2 + (#1)*\ax},{#3 + (#1)*\ay}}
	\providecommand{\rightangle}[3]{%
		\coordinate (ra1) at ($(#2)!2.4mm!(#1)$);
		\coordinate (ra3) at ($(#2)!2.4mm!(#3)$);
		\draw[line width=0.5pt] (ra1) -- ($(ra1)+(ra3)-(#2)$) -- (ra3);}
	\begin{minipage}[b]{0.50\textwidth}\centering
		\begin{tikzpicture}[>=Stealth, line width=0.8pt, scale=1.1]
			% proiezione: x in avanti (basso-sinistra), y destra, z su
			\def\ax{-0.65}\def\ay{-0.48}
			
			\def\az{0.7}                            % <1 = comprime in altezza
			\renewcommand{\xyz}[3]{{#2 + (#1)*\ax},{(#3)*\az + (#1)*\ay}}
			% componenti di P_k
			\def\xk{1.4}\def\yk{2.6}\def\zk{2.0}
			\coordinate (O)   at (\xyz{0}{0}{0});
			\coordinate (Px)  at (\xyz{\xk}{0}{0});
			\coordinate (Py)  at (\xyz{0}{\yk}{0});
			\coordinate (Pxy) at (\xyz{\xk}{\yk}{0});
			\coordinate (Pk)  at (\xyz{\xk}{\yk}{\zk});
			
			% --- scatola di decomposizione (base + verticale) ---
			\draw[gray!55, thin] (O) -- (Px) -- (Pxy) -- (Py) -- cycle;
			\draw[gray!55, thin] (Pxy) -- (Pk);
			
			% --- assi x,y,z ---
			\draw[->, thin] (O) -- (\xyz{2.6}{0}{0}) node[below left=-2pt] {$x$};
			\draw[->, thin] (O) -- (\xyz{0}{3.4}{0}) node[right]          {$y$};
			\draw[->, thin] (O) -- (\xyz{0}{0}{2.9}) node[above]          {$z$};
			
			% --- vettore posizione OP_k ---
			%				\draw[->, thin] (O) -- (Pk);
			%				\node[font=\small, above left=-8pt, yshift=6pt] at ($(O)!0.42!(Pk)$) {$\overrightarrow{OP_k}$};
			
			% --- terna delle componenti di forza in P_k ---
			\coordinate (Xt) at (\xyz{\xk+0.85}{\yk}{\zk});
			\coordinate (Yt) at (\xyz{\xk}{\yk+0.85}{\zk});
			\coordinate (Zt) at (\xyz{\xk}{\yk}{\zk+0.85});
			\draw[->, thick] (Pk) -- (Xt) node[above=4pt, xshift=-2pt] {$v_x(P)$};
			\draw[->, thick] (Pk) -- (Yt) node[right=-1pt] {$v_y(P)$};
			\draw[->, thick] (Pk) -- (Zt) node[above=-1pt] {$v_z(P)$};
			
			% --- etichette componenti (spigoli O->Py->Pxy->Pk) ---
			\node[font=\small, below=20pt, xshift=-20pt]      at ($(O)!0.5!(Py)$)   {$y(P)$};
			\node[font=\small,  right=1pt, yshift=-6pt] at ($(Py)!0.5!(Pxy)$) {$x(P)$};
			\node[font=\small, right, yshift=6pt]      at ($(Pxy)!0.5!(Pk)$) {$z(P)$};
			
			% --- punti ---
			\fill (O)  circle (1.5pt) node[above left=-1pt]  {$O$};
			\fill (Pk) circle (1.5pt) node[above right=-1pt] {$P$};
			
			
			% ===== MOMENTI in O (interni, assiali ->> + ricciolo) =====
			\coordinate (Mx) at (\xyz{1.2}{0}{0});
			\coordinate (My) at (\xyz{0}{1.2}{0});
			\coordinate (Mz) at (\xyz{0}{0}{1.4});
			\draw[->>, thick] (O) -- (Mx);
			%\draw[->, thin, rotate around={126:(Mx)}]([shift={(0.16,0)}]Mx) arc[x radius=0.16, y radius=0.06, start angle=0, end angle=305];
			\node[font=\small, above,  xshift=-4pt, yshift=6pt] at (Mx) {$\bm{\upmu}_{x}(O)$};
			\draw[->>, thick] (O) -- (My);
			%\draw[->, thin, rotate around={-90:(My)}] ([shift={(0.16,0)}]My) arc[x radius=0.16, y radius=0.06, start angle=0, end angle=305];
			\node[font=\small, below=-1pt, xshift=-6pt] at (My) {$\bm{\upmu}_{y}(O)$};
			\draw[->>, thick] (O) -- (Mz);
			%\draw[->, thin, rotate around={0:(Mz)}]	([shift={(0.16,0)}]Mz) arc[x radius=0.16, y radius=0.06, start angle=0, end angle=305];
			\node[font=\small, left=1pt] at (Mz) {$\bm{\upmu}_{z}(O)$};
			
			
		\end{tikzpicture}
	\end{minipage}
	\caption{Componenti del momento polare di $(P,\vb)$ rispetto al sistema cartesiano con origine nel polo $O$.}
	\label{f:componenti_momento_polare}
\end{figure}


\begin{oss}
	Il momento polare dipende dal polo $O$ ma non dal punto di
	applicazione $P$, purché questo rimanga sulla stessa linea
	d'azione. Infatti, se $P'$ giace sulla linea d'azione di $\vb$,
	allora $\overrightarrow{PP'} = \lambda\vb$ per qualche scalare
	$\lambda$, e quindi
	\begin{equation}
		\overrightarrow{OP'} \times \vb
		= (\overrightarrow{OP} + \lambda\vb) \times \vb
		= \overrightarrow{OP} \times \vb,
	\end{equation}
	essendo $\vb \times \vb = \zerovb$. Questa \textit{invarianza per
		traslazione lungo la linea d'azione} giustifica la notazione
	$\bm{\upmu}(O)$, che riporta il polo ma non il punto di
	applicazione (\FIGref{f:invarianza_retta_azione}).
	
	\begin{figure}[!htbp]
		\centering
		\def\ax{0.5}\def\ay{0.30}
		\providecommand{\xyz}[3]{{#2 + (#1)*\ax},{#3 + (#1)*\ay}}
		\providecommand{\rightangle}[3]{%
			\coordinate (ra1) at ($(#2)!2.4mm!(#1)$);
			\coordinate (ra3) at ($(#2)!2.4mm!(#3)$);
			\draw[line width=0.5pt] (ra1) -- ($(ra1)+(ra3)-(#2)$) -- (ra3);}
		\begin{tikzpicture}[>=Stealth, line width=0.8pt, scale=1.15]
			\draw[gray!60] (\xyz{-1.7}{-1.2}{0}) -- (\xyz{-1.7}{3.90}{0})-- (\xyz{4.7}{3.90}{0}) -- (\xyz{4.7}{-1.2}{0}) -- cycle;
			\node[gray!75] at (\xyz{-1.2}{-1.0}{0}) {$\pi$};
			
			\coordinate (O)   at (\xyz{0}{0}{0});
			\coordinate (Of)  at (\xyz{1.2}{0.3}{0});      % O + v   (base comune)
			\coordinate (Pk)  at (\xyz{0.5}{2.0}{0});      % P            (lambda = 0)
			\coordinate (Ff)  at (\xyz{1.7}{2.3}{0});      % P + v
			\coordinate (Pp)  at (\xyz{2.54}{2.51}{0});    % P'           (lambda = 1.7)
			\coordinate (Ffp) at (\xyz{3.74}{2.81}{0});    % P' + v
			\coordinate (La)  at (\xyz{0.14}{1.91}{0});    % retta d'azione, estremo a
			\coordinate (Lb)  at (\xyz{4.04}{2.89}{0});    % estremo b
			\coordinate (LOa) at (\xyz{-0.48}{-0.12}{0});  % retta per O // v, estremo a
			\coordinate (LOb) at (\xyz{1.92}{0.48}{0});    % estremo b
			\coordinate (MuB) at (\xyz{0}{0}{1.25});
			\coordinate (MuC) at (\xyz{0}{0}{1.47});
			
			% --- le due rette parallele della striscia (tratteggiate, dietro) ---
			%		\draw[gray!55, thin, dashed] (La) -- (Lb);    % retta d'azione di v
			%		\draw[gray!55, thin, dashed] (LOa) -- (LOb);  % retta per O parallela a v
			\draw[gray!55, thin, dashed] ($(O)!-0.6!(Of)$) -- ($(Of)!-1.0!(O)$);
			
			% --- i due parallelogrammi (entrambi grigi): base comune O-Of,
			%     lato opposto sulla stessa retta d'azione => stessa area ---
			\filldraw[gray!9, draw=gray!55, thin] (O) -- (Of) -- (Ffp) -- (Pp) -- cycle; % su OP'
			\filldraw[gray!9, draw=gray!55, thin] (O) -- (Of) -- (Ff)  -- (Pk) -- cycle; % su OP
			
			
			% --- OP e OP' ---
			\draw[->, thin] (O) -- (Pk);
			\node[font=\small, below right] at ($(O)!0.55!(Pk)$) {$\overrightarrow{OP}$};
			\draw[->, thin] (O) -- (Pp);
			\node[font=\small, above left=-4pt] at ($(O)!0.70!(Pp)$) {$\overrightarrow{OP'}$};
			
			% --- v applicato in P e in P' ---
			\draw[->, thick] (Pk) -- (Ff)  node[below right=-3pt] {$\vb$};
			\draw[->, thick] (Pp) -- (Ffp) node[right=-2pt] {$\vb$};
			
			% --- momento mu(O) ---
			\draw[->>, thick] (O) -- (MuB) node[left=1pt, yshift=-10pt] {$\bm{\upmu}(O)$};
			
			\begin{scope}[gray!55]
				\rightangle{MuB}{O}{Of}
				\rightangle{MuB}{O}{Pk}
			\end{scope}
			
			% --- punti ---
			\fill (O)  circle (1.5pt) node[left=2pt] {$O$};
			\fill (Pk) circle (1.5pt) node[below=2pt] {$P$};
			\fill (Pp) circle (1.5pt) node[below right=-2pt] {$P'$};
		\end{tikzpicture}
		\caption{Invarianza del momento polare per traslazione del punto
			di applicazione lungo la linea d'azione.}
		\label{f:invarianza_retta_azione}
	\end{figure}
\end{oss}

\subsection{Legge di variazione del momento al variare del polo}\index{momento!variazione al variare del polo}

Al variare del polo, il momento polare obbedisce alla legge affine
\begin{equation}\label{eq:trasporto_momento}
	\bm{\upmu}(P) = \bm{\upmu}(O) + \vb \times \overrightarrow{OP},
\end{equation}
che si dimostra scrivendo $\overrightarrow{PQ} =
\overrightarrow{PO} + \overrightarrow{OQ}$ e sfruttando
l'antisimmetria del prodotto vettoriale.

\begin{oss}
	Il termine $\vb \times \overrightarrow{OP}$ è ortogonale a $\vb$,
	pertanto la componente di $\bm{\upmu}$ nella direzione di $\vb$ è
	indipendente dal polo:
	\begin{equation}
		\bm{\upmu}(P) \cdot \vb = \bm{\upmu}(O) \cdot \vb
		\qquad \forall\, O,\, P \in \Ec.
	\end{equation}
	Questo invariante scalare si ritroverà come invariante dei sistemi
	di forze e dei campi di spostamento rigido.
\end{oss}


\subsection{Momento assiale}\label{sec:momento_assiale}

Il momento polare porta con sé informazioni su tutte le direzioni dello
spazio. Quando interessa l'effetto rotazionale attorno a una singola
retta, si introduce il \textit{momento assiale}\index{momento!assiale|textbf}\nomenclature[G]{$\mu_r$}{momento assiale}. Dato un vettore
applicato $(P,\,\vb)$ e una retta $r$ di versore $\ab_r$ passante per
$O$, si definisce momento assiale di $\vb$ rispetto a $r$ lo scalare
\begin{equation}\label{eq:momento_assiale}
	\mu_r \coloneqq \bm{\upmu}(O) \cdot \ab_r
	= (\overrightarrow{OP} \times \vb) \cdot \ab_r.
\end{equation}
La \eqref{eq:momento_assiale} è un prodotto misto: $\mu_r$ è dunque
il volume con segno del parallelepipedo di spigoli
$\overrightarrow{OP}$, $\vb$ e $\ab_r$
(\FIGref{f:momento_assiale}a), positivo se la terna è destrorsa.
Questa lettura geometrica rende trasparenti le proprietà di
invarianza discusse di seguito.

Le componenti cartesiane $\mu_x(O)$, $\mu_y(O)$, $\mu_z(O)$ del
momento polare sono i momenti assiali rispetto ai tre assi coordinati
passanti per $O$.


%%%===================================================================================================
\begin{figure}[!htbp]
	\centering
	\begin{minipage}[b]{0.44\linewidth}
		\centering
		\begin{tikzpicture}[>=Stealth,
			x={(0.85cm,0cm)}, y={(0cm,0.85cm)}, z={(0.425cm,0.255cm)},% proiezione obliqua (ax=0.5, ay=0.30) a scala 0.85
			vec/.style={->,thick},
			edge/.style={gray!60,line width=0.6pt},
			pt/.style={circle,fill,inner sep=1.1pt}]
			% ---- vertici: parallelepipedo di spigoli OP, v, a_r (vertice in O) ----
			\coordinate (O)   at (0, 0, 0);
			\coordinate (Ar)  at (-0.5, 1.0, -0.4);   % O + a_r
			\coordinate (P)   at (3.0, 0.2, 1.2);     % O + OP
			\coordinate (V0)  at (1.0, 1.1, -0.8);    % O + v
			\coordinate (PV)  at (4.0, 1.3, 0.4);     % P + v
			\coordinate (PA)  at (2.5, 1.2, 0.8);     % P + a_r
			\coordinate (PVA) at (3.5, 2.3, 0.0);     % P + v + a_r
			\coordinate (VA0) at (0.5, 2.1, -1.2);    % O + v + a_r
			% ---- retta r ----
			\coordinate (Re) at (-1.2, 2.4, -0.96);       % estremo alto di r
			\draw[line width=0.6pt] (0.45, -0.9, 0.36) -- (Re)
			node[above=1pt] {$r$};
			% ---- spigoli sottili del parallelepipedo ----
			\draw[edge] (O) -- (V0) -- (PV);
			\draw[edge] (Ar) -- (PA) -- (PVA) -- (VA0) -- cycle;
			\draw[edge] (P) -- (PA);
			\draw[edge] (PV) -- (PVA);
			\draw[edge] (V0) -- (VA0);
			% ---- momento polare mu(O) e sua proiezione mu_r su r ----
			\coordinate (Mu) at (-1.07, 2.60, 2.24);      % O + mu(O), riscalato
			\coordinate (Fr) at (-0.794, 1.589, -0.636);  % piede della proiezione su r
			\draw[gray!55,thin,dashed] (Mu) -- (Fr);
			\coordinate (ra1) at ($(Fr)!2.4mm!(Re)$);
			\coordinate (ra3) at ($(Fr)!2.4mm!(Mu)$);
			\draw[line width=0.5pt] (ra1) -- ($(ra1)+(ra3)-(Fr)$) -- (ra3);
			% quota di mu_r lungo r (stanghette lunghe per scavalcare l'etichetta del versore)
			\draw[gray!55,thin] (O)  -- ($(O)!1.05cm!90:(Fr)$);
			\draw[gray!55,thin] (Fr) -- ($(Fr)!1.05cm!-90:(O)$);
			\draw[<->,thin] ($(O)!0.88cm!90:(Fr)$) -- ($(Fr)!0.88cm!-90:(O)$)
			node[midway,anchor=east,xshift=-2pt,font=\small] {$\mu_r$};
			\draw[->>,thick] (O) -- (Mu)
			node[pos=0.80,anchor=west,xshift=3pt] {$\bm{\upmu}(O)$};
			% ---- vettori (spessi): a_r, OP, v ----
			\draw[vec] (O) -- (Ar) node[midway,anchor=east,xshift=-2pt] {$\ab_r$};
			\draw[vec] (O) -- (P)  node[midway,below=2pt] {$\overrightarrow{OP}$};
			\draw[vec] (P) -- (PV) node[midway,anchor=west,xshift=3pt] {$\vb$};
			% ---- punti ----
			\node[pt] at (O) {};
			\node[below=3pt] at (O) {$O$};
			\node[pt] at (P) {};
			\node[below right=1pt] at (P) {$P$};
		\end{tikzpicture}\\[4pt]
		{\small (a)}
	\end{minipage}
	\hfill
	\begin{minipage}[b]{0.54\linewidth}
		\centering
		\begin{tikzpicture}[>=Stealth,
			x={(0.85cm,0cm)}, y={(0cm,0.85cm)}, z={(0.425cm,0.255cm)},% proiezione obliqua (ax=0.5, ay=0.30) a scala 0.85
			vec/.style={->,thick},
			edge/.style={gray!60,line width=0.6pt},
			edgeb/.style={gray!60,line width=0.6pt,dashed},
			pt/.style={circle,fill,inner sep=1.1pt}]
			% ---- vertici comuni (come nel pannello a) ----
			\coordinate (O)   at (0, 0, 0);
			\coordinate (Ar)  at (-0.5, 1.0, -0.4);   % O + a_r
			\coordinate (P)   at (3.0, 0.2, 1.2);     % O + OP
			\coordinate (V0)  at (1.0, 1.1, -0.8);    % O + v
			\coordinate (PV)  at (4.0, 1.3, 0.4);     % P + v
			\coordinate (PA)  at (2.5, 1.2, 0.8);     % P + a_r
			\coordinate (PVA) at (3.5, 2.3, 0.0);     % P + v + a_r
			\coordinate (VA0) at (0.5, 2.1, -1.2);    % O + v + a_r
			% ---- vertici del secondo parallelepipedo: O' = O + 1.6 a_r ----
			\coordinate (Op)   at (-0.8, 1.6, -0.64); % O'
			\coordinate (OpV)  at (0.2, 2.7, -1.44);  % O' + v
			\coordinate (OpA)  at (-1.3, 2.6, -1.04); % O' + a_r
			\coordinate (OpVA) at (-0.3, 3.7, -1.84); % O' + v + a_r
			% ---- retta r (prolungata oltre O') ----
			\draw[line width=0.6pt] (0.45, -0.9, 0.36) -- (-1.5, 3.0, -1.2)
			node[above=1pt] {$r$};
			% ---- parallelepipedo 1 (vertice in O): spigoli sottili pieni ----
			\draw[edge] (O) -- (V0) -- (PV);
			\draw[edge] (Ar) -- (PA);
			\draw[edge] (V0) -- (VA0);
			\draw[edge] (Ar) -- (VA0);
			\draw[edge] (VA0) -- (PVA);
			% ---- faccia comune in P (condivisa dai due solidi) ----
			\draw[edge] (PV) -- (PVA) -- (PA) -- (P);
			% ---- parallelepipedo 2 (vertice in O'): spigoli tratteggiati ----
			\draw[edgeb] (Op) -- (OpV) -- (OpVA) -- (OpA) -- cycle;
			\draw[edgeb] (OpV) -- (PV);
			\draw[edgeb] (OpA) -- (PA);
			\draw[edgeb] (OpVA) -- (PVA);
			% ---- vettori (spessi) ----
			\draw[vec] (O) -- (Ar) node[midway,anchor=east,xshift=-2pt] {$\ab_r$};
			\draw[vec] (O) -- (P)  node[midway,below=2pt] {$\overrightarrow{OP}$};
			\draw[vec] (Op) -- (P) node[pos=0.45,below] {$\overrightarrow{O'P}$};
			\draw[vec] (P) -- (PV) node[midway,anchor=west,xshift=3pt] {$\vb$};
			% ---- punti ----
			\node[pt] at (O) {};
			\node[below=3pt] at (O) {$O$};
			\node[pt] at (Op) {};
			\node[anchor=east,xshift=-3pt] at (Op) {$O'$};
			\node[pt] at (P) {};
			\node[below right=1pt] at (P) {$P$};
		\end{tikzpicture}\\[4pt]
		{\small (b)}
	\end{minipage}
	\caption[Momento assiale di un vettore applicato]{Momento assiale
		di un vettore applicato.
		\textit{(a)}~Il momento assiale è la proiezione $\mu_r$ del momento
		polare $\bm{\upmu}(O)$ sulla retta $r$ e coincide con il volume con
		segno del parallelepipedo di spigoli $\protect\overrightarrow{OP}$,
		$\vb$ e $\ab_r$.
		\textit{(b)}~Invarianza rispetto al punto scelto sulla retta: i
		parallelepipedi costruiti da $O$ e da $O'$ condividono la faccia
		in $P$ generata da $\vb$ e $\ab_r$, e le facce opposte giacciono
		entrambe nel piano per $r$ parallelo a $\vb$; i volumi coincidono
		per scorrimento.}
	\label{f:momento_assiale}
\end{figure}
%%%===================================================================================================


Il momento assiale gode di una \textit{doppia invarianza}: non dipende
dal punto di applicazione $P$ (per la stessa ragione del momento
polare), né dal punto $O$ scelto sulla retta $r$. Quest'ultima
proprietà si dimostra osservando che, se $O' = O + \lambda\ab_r$,
allora
\begin{equation}
	(\overrightarrow{O'P} \times \vb) \cdot \ab_r
	= [(\overrightarrow{OP} - \lambda\ab_r) \times \vb] \cdot \ab_r
	= (\overrightarrow{OP} \times \vb) \cdot \ab_r,
\end{equation}
essendo $(\ab_r \times \vb) \cdot \ab_r = 0$. Questa doppia invarianza
giustifica la notazione $\mu_r$, che riporta solo la retta di
riferimento.

La dimostrazione algebrica ha una controparte geometrica immediata:
i parallelepipedi costruiti da $O$ e da $O'$ condividono la faccia
in $P$ generata da $\vb$ e $\ab_r$, mentre le facce opposte giacciono
entrambe nel piano per $r$ parallelo a $\vb$; i due volumi coincidono
quindi per scorrimento (\FIGref{f:momento_assiale}b), così come, nel
piano, coincidevano le aree dei parallelogrammi al variare del punto
di applicazione lungo la retta d'azione
(\FIGref{f:invarianza_retta_azione}). Lo stesso argomento, con
scorrimento lungo $\vb$, rilegge l'invarianza rispetto al punto di
applicazione~$P$.


La doppia invarianza ha una conseguenza operativa: il calcolo del
momento assiale si riduce a un problema piano
(\FIGref{f:momento_assiale_piano}). Si scelga come piano di
riferimento il piano $\pi$ passante per il punto di applicazione
$P$ e ortogonale alla retta, e come polo l'intersezione
$O = r \cap \pi$. Decomposto il vettore rispetto alla retta,
$\vb = \vb_\parallel + \vb_\perp$, con $\vb_\parallel$ parallelo ad
$\ab_r$ e $\vb_\perp$ giacente in $\pi$, la componente parallela
non contribuisce: il parallelepipedo di spigoli
$\overrightarrow{OP}$, $\vb_\parallel$ e $\ab_r$ è degenere, avendo
due spigoli paralleli. Resta dunque
\begin{equation}
	\mu_r = (\overrightarrow{OP} \times \vb_\perp) \cdot \ab_r
	= |\vb_\perp|\,d,
\end{equation}
dove $d$ è la distanza della linea d'azione di $\vb_\perp$ dalla
retta --- il braccio $\overline{OH}$ della
\FIGref{f:momento_assiale_piano} --- e il segno è positivo se la
rotazione indotta da $\vb_\perp$ attorno a $r$ è concorde con
$\ab_r$ secondo la regola della mano destra, ovvero antioraria per
chi guarda $\pi$ dalla punta di $\ab_r$. Il momento assiale è
quindi il momento piano della componente ortogonale rispetto al
polo $O$. In particolare, $\mu_r$ si annulla se e solo se il
parallelepipedo della \FIGref{f:momento_assiale}\,(a)
degenera, cioè se e solo se la retta $r$ e la linea d'azione di
$\vb$ sono complanari: incidenti ($d = 0$) o parallele
($\vb_\perp = \zerovb$).

%%%===================================================================================================
\begin{figure}[htbp]
	\centering
	\begin{tikzpicture}[>=Stealth,
		x={(0.85cm,0cm)}, y={(0cm,0.85cm)}, z={(0.425cm,0.255cm)},% proiezione obliqua (ax=0.5, ay=0.30) a scala 0.85
		vec/.style={->,thick},
		aux/.style={->,thin,gray!60},
		lact/.style={gray!55,thin,dashed},
		pt/.style={circle,fill,inner sep=1.1pt}]
		% ---- punti principali (coordinate: orizz, vert, profondita' [recede in alto a dx]) ----
		\coordinate (O)   at (0, 0, 0);            % O = r intersezione sigma
		\coordinate (P)   at (2.6, 0, 0.9);        % punto di applicazione, NEL piano
		\coordinate (Pq)  at (3.1, 0, 2.0);        % P + v_perp (nel piano)
		\coordinate (Pv)  at (3.1, 1.0, 2.0);      % P + v
		\coordinate (Hf)  at (1.816, 0, -0.825);   % piede del braccio d
		\coordinate (La)  at (1.55, 0, -1.41);     % linea d'azione di v_perp, estremo a
		\coordinate (Lb)  at (3.175, 0, 2.165);    % estremo b
		\coordinate (Ra)  at (0, 1.1, 0);          % O + a_r
		\coordinate (Re)  at (0, 2.6, 0);          % estremo alto di r
		% ---- piano \pi per P ortogonale a r ----
		\draw[gray!60] (-2.1, 0, -2.6) -- (4.5, 0, -2.6) -- (4.5, 0, 2.8) -- (-2.1, 0, 2.8) -- cycle;
		%\draw[gray!60] (-2.1, 0, -2.1) -- (4.5, 0, -2.1) -- (4.5, 0, 2.8)	-- (-2.1, 0, 2.8) -- cycle;
		\node[gray!75] at (-1.6, 0, -1.8) {$\pi$};
		% ---- retta r (perfora il piano in O) e versore ----
		\draw[line width=0.6pt] (0, -1.0, 0) -- (Re) node[above=1pt] {$r$};
		\coordinate (raA1) at ($(O)!2.4mm!(Re)$);
		\coordinate (raA3) at ($(O)!2.4mm!(Hf)$);
		\draw[line width=0.5pt] (raA1) -- ($(raA1)+(raA3)-(O)$) -- (raA3);
		\draw[vec] (O) -- (Ra) node[pos=0.95,anchor=east,xshift=-3pt] {$\ab_r$};
		
		% ---- verso positivo di rotazione attorno a r (mano destra, antiorario da sopra) ----
		\draw[thin, postaction={decorate},
		decoration={markings, mark=at position 1 with {\arrow{Stealth}}}]
		plot[domain=15:255,samples=90] ({0.52*cos(\x)}, 0.5, {0.52*sin(\x)});
		
		
		% ---- linea d'azione di v_perp nel piano e braccio d ----
		\draw[lact] (La) -- (Lb);
		\draw[thin] (O) -- (Hf) node[midway,anchor=south,yshift=2pt,font=\small] {$d$};
		\coordinate (raB1) at ($(Hf)!2.4mm!(Lb)$);
		\coordinate (raB3) at ($(Hf)!2.4mm!(O)$);
		\draw[line width=0.5pt] (raB1) -- ($(raB1)+(raB3)-(Hf)$) -- (raB3);
		% ---- decomposizione di v in P: v_perp nel piano, v_par lungo r ----
		\draw[aux] (Pq) -- (Pv) node[midway,anchor=west,xshift=2pt,gray!75,font=\footnotesize] {$\vb_\parallel$};
		\draw[vec] (P) -- (Pv) node[pos=0.55,anchor=south east,xshift=-1pt] {$\vb$};
		\draw[vec] (P) -- (Pq) node[pos=0.55,anchor=north west,xshift=2pt,yshift=-1pt] {$\vb_\perp$};
		% ---- punti ----
		\node[pt] at (O) {};
		\node[anchor=north east,xshift=-2pt] at (O) {$O$};
		\node[pt] at (P) {};
		\node[below right=1pt] at (P) {$P$};
		\node[pt] at (Hf) {};
		\node[anchor=north west,xshift=2pt,yshift=1pt] at (Hf) {$H$};
	\end{tikzpicture}
	\caption[Momento assiale come momento piano]{Riduzione al piano del
		momento assiale: per la doppia invarianza si sceglie il piano
		$\pi$ per $P$ ortogonale a $r$ e, come polo, l'intersezione
		$O = r \cap \pi$. Decomposto $\vb = \vb_\parallel + \vb_\perp$,
		la componente parallela non contribuisce e $\mu_r = |\vb_\perp|\,d$
		è il momento piano di $\vb_\perp$ rispetto a $O$, positivo se la
		rotazione indotta è concorde con $\ab_r$ secondo la regola della
		mano destra.}
	\label{f:momento_assiale_piano}
\end{figure}
%%%===================================================================================================



%-------------------------------------------------------------------------------------------------------------------
\section{Vettori polari e vettori assiali}\label{sec:polari_assiali}
%-------------------------------------------------------------------------------------------------------------------

L'algebra vettoriale produce due tipi di oggetti che si comportano in
modo diverso sotto riflessioni del sistema di riferimento. La
distinzione non è formale: essa ha conseguenze fisiche precise e sarà
sfruttata sistematicamente nell'analisi delle strutture.

\paragraph{Vettori polari.}\index{vettore!polare|textbf}
Sono vettori che agiscono nella direzione e nel verso da essi definiti.
Sotto una riflessione rispetto a un piano, la componente perpendicolare
al piano cambia segno, le componenti parallele restano invariate.
Sono vettori polari le forze, gli spostamenti, le velocità, i vettori
posizione.

\paragraph{Vettori assiali (o pseudovettori).}\index{vettore!assiale|textbf}\index{pseudovettore|see{vettore (assiale)}}
Sono vettori che si ottengono come prodotto vettoriale di due vettori
polari e che agiscono sul piano ortogonale alla loro direzione. Sotto
una riflessione, acquisiscono un segno aggiuntivo rispetto a un vettore
polare soggetto alla stessa trasformazione: la componente perpendicolare
al piano di riflessione resta invariata, le componenti parallele cambiano
segno. Sono vettori assiali il momento di una forza, la velocità
angolare, la coppia.


\subsection{Comportamento sotto riflessione}

Una riflessione rispetto a un piano $\pi$ di normale unitaria $\nb$ è
rappresentata dall'operatore
\begin{equation}
	\Rbb_{\nb}(\vb) \coloneqq \vb - 2(\vb \cdot \nb)\,\nb,
\end{equation}
che inverte la componente di $\vb$ parallela a $\nb$ e lascia invariata
quella perpendicolare (\FIGref{fig:riflessione}). La riflessione è una
trasformazione ortogonale impropria: $\Rbb_{\nb}^\top \Rbb_{\nb} = \Ib$
e $\det\Rbb_{\nb} = -1$.


%%===================================================================================================
\begin{figure}[!htbp]
	\centering
	% ---- helper locali (blocco autocontenuto) --------------------------------------------------
	\providecommand{\rightangle}[3]{\coordinate (ra1) at ($(#2)!2.2mm!(#1)$);
		\coordinate (ra3) at ($(#2)!2.2mm!(#3)$);
		\draw[gray!55,line width=0.5pt](ra1)--($(ra1)+(ra3)-(#2)$)--(ra3);}% angolo retto in #2
	\begin{tikzpicture}[>=Stealth,scale=1.0,
		vec/.style={->,thick},
		aux/.style={->,thin,gray!60},
		auxlab/.style={gray!75,font=\footnotesize},
		pt/.style={circle,fill,inner sep=1.1pt},
		mirror/.style={thin,gray!55,dashed}]
		%
		% ---- traccia del piano di riflessione (asse di simmetria) -------------------------------
		\coordinate (A) at (235:3.3);          % estremo inferiore
		\coordinate (B) at (55:3.3);           % estremo superiore
		\draw[very thick] (A) -- (B);
		\node[anchor=west] at ($(B)+(0.10,0.08)$) {$\uppi$};
		% simbolo dell'asse di simmetria: due tratti passanti per lo stesso punto, inclinati
		% di +-35 rispetto all'asse (che ne e' bisettrice); in punta a ciascuno un pennone
		% triangolare con il cateto superiore ortogonale all'asta e un cateto sul tratto stesso;
		% #1 = posizione del punto sull'asse
		\providecommand{\symaxis}[1]{%
			\draw[line width=0.6pt] ($(55:#1)+(200:0.40)$) -- ($(55:#1)+(20:0.55)$)
			-- ($(55:#1)+(20:0.55)+(-70:0.18)$) -- ($(55:#1)+(20:0.31)$);
			\draw[line width=0.6pt] ($(55:#1)+(270:0.40)$) -- ($(55:#1)+(90:0.55)$)
			-- ($(55:#1)+(90:0.55)+(180:0.18)$) -- ($(55:#1)+(90:0.31)$);}
		\symaxis{-2.85}
		% normale al piano
		\coordinate (N0)   at (55:2.85);
		\coordinate (Ntip) at ($(N0)+(-35:0.85)$);
		\draw[vec] (N0) -- (Ntip) node[anchor=west,xshift=1pt] {$\nb$};
		\rightangle{B}{N0}{Ntip}
		%
		%==========================================================================================
		% COPPIA SUPERIORE: vettore polare
		%==========================================================================================
		\coordinate (F1) at (55:1.2);                       % piede sulla traccia
		\coordinate (P)  at ($(F1)+(-35:1.9)$);             % punto originale
		\coordinate (Ps) at ($(F1)+(145:1.9)$);             % punto simmetrico
		\draw[mirror] (P) -- (Ps);
		% --- lato originale: v = comp.parallela + comp.normale
		\coordinate (Pu) at ($(P)+(55:1.6)$);               % P + (v - (v.n)n)
		\coordinate (Pv) at ($(Pu)+(-35:0.7)$);             % P + v
		\draw[aux] (P) -- (Pu) node[auxlab,pos=0.6,sloped,above=1pt] {$\vb-(\vb\cdot\nb)\,\nb$};
		\draw[aux] (Pu) -- (Pv);
		\node[auxlab,anchor=west] at ($(Pu)!0.5!(Pv)+(0.10,0.05)$) {$(\vb\cdot\nb)\,\nb$};
		\draw[vec] (P) -- (Pv) node[midway,below=2pt,yshift=-1pt] {$\vb$};
		\node[pt] at (P) {};
		\node[below=3pt] at (P) {$P$};
		% --- lato riflesso: R_n v = comp.parallela - comp.normale
		\coordinate (Qu) at ($(Ps)+(55:1.6)$);              % Ps + (v - (v.n)n)
		\coordinate (Qv) at ($(Qu)+(145:0.7)$);             % Ps + R_n v
		\draw[aux] (Ps) -- (Qu) node[auxlab,pos=0.55,sloped,below=1pt] {$\vb-(\vb\cdot\nb)\,\nb$};
		\draw[aux] (Qu) -- (Qv);
		\node[auxlab,anchor=west] at ($(Qu)!0.5!(Qv)+(0.12,0.10)$) {$-(\vb\cdot\nb)\,\nb$};
		\draw[vec] (Ps) -- (Qv) node[midway,anchor=east,xshift=-3pt] {$\Rbb_{\nb}\vb$};
		\node[pt] at (Ps) {};
		\node[anchor=east,xshift=-3pt,yshift=-2pt] at (Ps) {$P_\star$};
		%
		%==========================================================================================
		% COPPIA INFERIORE: vettore assiale
		%==========================================================================================
		\coordinate (F2) at (235:1.4);                      % piede sulla traccia
		\coordinate (Q)  at ($(F2)+(-35:1.7)$);             % punto originale
		\coordinate (Qs) at ($(F2)+(145:1.7)$);             % punto simmetrico
		\draw[mirror] (Q) -- (Qs);
		% --- lato originale: v, w e parallelogramma; circolazione antioraria (omega uscente)
		\coordinate (Wv)  at ($(Q)+(55:0.5)+(-35:1.2)$);    % Q + v
		\coordinate (Ww)  at ($(Q)+(55:1.4)+(145:0.2)$);    % Q + w
		\coordinate (Wvw) at ($(Wv)+(55:1.4)+(145:0.2)$);   % Q + v + w
		\draw[thin,gray!60] (Wv) -- (Wvw) -- (Ww);
		\node[auxlab] at ($(Q)!0.5!(Wvw)$) {$|\bm{\upomega}|$};
		\draw[vec] (Q) -- (Wv) node[midway,below=2pt,xshift=2pt] {$\vb$};
		\draw[vec] (Q) -- (Ww) node[midway,anchor=east,xshift=-2pt] {$\wb$};
		\draw[thick,->] ($(Q)+(60:0.30)$) arc[start angle=60,end angle=330,radius=0.30];
		\node[pt] at (Q) {};
		\node[anchor=east,xshift=-9pt,yshift=4pt]  at (Q) {$P$};
		\node[anchor=north east,xshift=-2pt,yshift=-9pt] at (Q) {$\bm{\upomega}$};
		\node at ($(Q)+(0.85,-1.05)$) {$\bm{\upomega}=\vb\times\wb$};
		% --- lato riflesso: R_n v, R_n w; circolazione oraria (omega_star entrante)
		\coordinate (Zv)  at ($(Qs)+(55:0.5)+(145:1.2)$);   % Qs + R_n v
		\coordinate (Zw)  at ($(Qs)+(55:1.4)+(-35:0.2)$);   % Qs + R_n w
		\coordinate (Zvw) at ($(Zv)+(55:1.4)+(-35:0.2)$);   % Qs + R_n v + R_n w
		\draw[thin,gray!60] (Zv) -- (Zvw) -- (Zw);
		\node[auxlab] at ($(Qs)!0.5!(Zvw)$) {$|\bm{\upomega}_\star|$};
		\draw[vec] (Qs) -- (Zv) node[midway,anchor=east,xshift=-3pt] {$\Rbb_{\nb}\vb$};
		\draw[vec] (Qs) -- (Zw) node[midway,anchor=west,xshift=4pt,yshift=-3pt] {$\Rbb_{\nb}\wb$};
		\draw[thick,->] ($(Qs)+(120:0.30)$) arc[start angle=120,end angle=-150,radius=0.30];
		\node[pt] at (Qs) {};
		\node[anchor=east,xshift=-9pt,yshift=2pt] at (Qs) {$P_\star$};
		\node[anchor=north,yshift=-9pt] at (Qs) {$\bm{\upomega}_\star$};
		\node at ($(Qs)+(-0.15,-1.15)$) {$\bm{\upomega}_\star=-\Rbb_{\nb}\bm{\upomega}$};
	\end{tikzpicture}
	\caption{Riflessione rispetto al piano $\uppi$ di normale $\nb$. In alto:
		il vettore polare $\vb$, applicato in $P$, si trasforma nel simmetrico
		$\Rbb_{\nb}\vb$, applicato nel punto simmetrico $P_\star$: la componente
		parallela al piano resta invariata, quella normale si inverte. In basso:
		i riflessi $\Rbb_{\nb}\vb$ e $\Rbb_{\nb}\wb$ definiscono una circolazione
		di verso opposto, e il vettore assiale $\bm{\upomega}=\vb\times\wb$ si
		trasforma nell'antisimmetrico
		$\bm{\upomega}_\star=-\Rbb_{\nb}(\bm{\upomega})$.}
	\label{fig:riflessione}
\end{figure}
%%%===================================================================================================


Un vettore polare $\vb(P)$ si trasforma sotto riflessione nel vettore	simmetrico $\Rbb_{\nb}(\vb(P))$, applicato nel punto simmetrico	$P_\star$ (\FIGref{fig:riflessione}). Un vettore assiale $\bm{\upomega} = \ub \times \vb$ si trasforma
invece nel vettore antisimmetrico. Applicando la riflessione separatamente a $\ub$ e $\vb$, i vettori riflessi $\Rbb_{\nb}\ub$
e $\Rbb_{\nb}\vb$ giacciono nel piano riflesso; la regola della mano destra applicata alla rotazione da $\Rbb_{\nb}\ub$ verso
$\Rbb_{\nb}\vb$ attraverso l'angolo convesso fornisce però un verso	opposto a quello di $\Rbb_{\nb}\bm{\upomega}$, perché la riflessione	ha invertito l'orientazione dello spazio. Il vettore assiale si	trasforma quindi nel vettore antisimmetrico:
\begin{equation}\label{eq:rif_assiale}
	\bm{\upomega}_\star(P_\star) = -\Rbb_{\nb}(\bm{\upomega}(P)).
\end{equation}



\subsection{Regole di composizione}

Il comportamento sotto riflessione determina le seguenti regole,
verificabili direttamente dalla~\eqref{eq:rif_assiale}:
\begin{enumerate}[label=(\roman*)]
	\item il prodotto vettoriale di due vettori \textit{polari} è un
	vettore \textit{assiale};
	\item il prodotto vettoriale di un vettore \textit{polare} e uno
	\textit{assiale} è un vettore \textit{polare};
	\item il prodotto vettoriale di due vettori \textit{assiali} è un
	vettore \textit{assiale}.
\end{enumerate}
Il momento di una forza, ad esempio, è il prodotto vettoriale di un
vettore posizione (polare) per una forza (polare): è dunque un vettore
assiale, coerentemente con il suo significato fisico di rotazione
attorno a un asse.

\begin{oss}
	Nel passaggio da una base destra a una base sinistra, ottenuta	invertendo tutti e tre i versori ($\ab_i \to -\ab_i$), le
	componenti dei vettori polari cambiano tutte segno	($u_i \to -u_i$), mentre quelle dei vettori assiali restano
	invariate ($\omega_i \to \omega_i$). Le regole di composizione	seguono direttamente da questa oss: nel prodotto
	vettoriale $(\ub \times \vb)_i = \varepsilon_{ijk}\,u_j\,v_k$,	il segno del risultato è determinato dal prodotto dei segni delle
	componenti dei fattori. Pertanto polare$\times$polare produce un	assiale (i segni si elidono a coppie), polare$\times$assiale
	produce un polare (un solo cambio di segno), assiale$\times$assiale	produce un assiale (nessun cambio di segno).
\end{oss}

\subsection{Rappresentazione grafica}

Nei disegni strutturali i due tipi di vettori si distinguono mediante
simboli diversi. I vettori polari sono rappresentati con una
\textit{freccia semplice}: il segmento indica la linea d'azione, la
punta indica il verso. I vettori assiali sono rappresentati con una
\textit{doppia freccia} (o con una freccia semicircolare nel piano
ortogonale): essi non agiscono lungo la propria direzione, ma inducono
una rotazione nel piano ad essa ortogonale, il cui verso è determinato
dalla regola della mano destra.

\begin{center}
	\begin{tikzpicture}[>=Stealth, line width=0.8pt]
		% ── Vettore polare ──
		\begin{scope}[xshift=0cm]
			\fill (0,0) circle (1.5pt) node[below=3pt]{$P$};
			\draw[->, thick] (0,0) -- (2.5,0)
			node[midway, above=2pt]{$\fb$};
			\node[below] at (1.25,-0.6) {\small vettore polare};
		\end{scope}
		% ── Vettore assiale ──
		\begin{scope}[xshift=6.5cm]
			% piano ortogonale al vettore centrato in P
			\draw[gray!50, fill=gray!10]
			(0,0) ellipse [x radius=0.2, y radius=0.7];
			% freccia semicircolare nel piano dell'ellisse
			% applicando la stessa trasformazione dell'ellisse
			\draw[<-, thin, xscale=0.2/0.7]
			(-0.55, 0) arc[start angle=180,
			end angle=0, radius=0.55];
			% vettore assiale
			\fill (0,0) circle (1.5pt) node[below=3pt]{$P$};
			\draw[->>, thick] (0,0) -- (2.5,0)
			node[midway, above=2pt]{$\bm{\upmu}$};
			\node[below] at (1.25,-0.6) {\small vettore assiale};
		\end{scope}
	\end{tikzpicture}
\end{center}

\noindent Una forza (polare) spinge il corpo nella direzione della
freccia; un momento (assiale) tende a far ruotare il corpo nel piano
ortogonale alla doppia freccia.


%-------------------------------------------------------------------------------------------------------------------
\section{Il caso piano}\label{sec:caso_piano}
%-------------------------------------------------------------------------------------------------------------------

Nell'analisi delle strutture piane i punti appartengono al piano
$(x,y)$ e le grandezze vettoriali si semplificano notevolmente.
È utile richiamare esplicitamente queste semplificazioni, poiché
costituiscono il riferimento costante nel seguito del testo.

\paragraph{Vettori polari nel piano.}
Un vettore polare $\vb$ giacente nel piano $(x,y)$ ha due sole
componenti non nulle:
\begin{equation}
	\vb = v_x\,\ab_x + v_y\,\ab_y, \qquad
	\bm{v} = \begin{Bmatrix} v_x \\ v_y \end{Bmatrix} \in \mathbb{R}^2.
\end{equation}

\paragraph{Vettori assiali nel piano.}
Un vettore assiale generato da due vettori del piano è necessariamente
ortogonale al piano stesso: ha come unica componente non nulla quella
lungo $\ab_z$. Il momento polare di un vettore applicato $(P,\,\vb)$
rispetto a un polo $O$ nel piano si riduce dunque a uno
\textit{scalare con segno}:
\begin{equation}\label{eq:momento_piano}
	\mu_z(O) = x(P)\,v_y - y(P)\,v_x,
\end{equation}
dove $x(P)$ e $y(P)$ sono le coordinate di $P$ rispetto a $O$.
Il segno è positivo se la rotazione indotta è antioraria, negativo
se oraria.

\paragraph{Rappresentazione grafica nel piano.}
I vettori polari si rappresentano con la consueta freccia nel piano.
I vettori assiali, essendo ortogonali al piano del disegno, si
rappresentano in due modi alternativi: con una freccia semicircolare
che indica direttamente il verso di rotazione, oppure con un simbolo
puntuale --- il cerchio con un punto ($\odot$) per un vettore uscente
dal piano, il cerchio con una croce ($\otimes$) per un vettore
entrante. Nel  testo si farà uso prevalentemente della
freccia semicircolare.


\begin{center}
	\begin{tikzpicture}[>=Stealth, line width=0.8pt]
		% ── Forza (polare, nel piano) ──
		\begin{scope}[xshift=0cm]
			\fill (0,0) circle (1.5pt) node[below left=2pt]{$P$};
			\draw[->, thick] (0,0) -- (2,0.8)
			node[above left=1pt]{$\fb$};
			\node[below] at (1,-0.6) {\small forza (polare)};
		\end{scope}
		% ── Momento uscente ──
		\begin{scope}[xshift=5cm]
			% freccia semicircolare antioraria
			\draw[->] (0.4,0) arc[start angle=0,
			end angle=180, radius=0.4];
			\draw (0,0) circle (0.2);
			\fill (0,0) circle (1.5pt);
			\node[below=6pt] at (0,0) {$\mu_z > 0$};
			\node[below] at (0,-0.6) {\small momento uscente};
		\end{scope}
		% ── Momento entrante ──
		\begin{scope}[xshift=9cm]
			% freccia semicircolare oraria
			\draw[->] (-0.4,0) arc[start angle=180,
			end angle=0, radius=0.4];
			\draw (0,0) circle (0.2);
			\draw (-0.12,-0.12) -- (0.12,0.12);
			\draw (-0.12,0.12) -- (0.12,-0.12);
			\node[below=6pt] at (0,0) {$\mu_z < 0$};
			\node[below] at (0,-0.6) {\small momento entrante};
		\end{scope}
	\end{tikzpicture}
\end{center}


\section{Esercizi proposti}

\begin{esercizio}[Prodotto scalare e prodotto vettoriale di due vettori]
	\label{ex:A03_prodotti}
	
	Si considerino i due vettori
	\begin{equation*}
		\bm{a} = (2\;\,1\;\,-1)^\top, \qquad
		\bm{b} = (1\;\,3\;\,2)^\top.
	\end{equation*}
	(a) Calcolare il prodotto scalare $\ab \cdot \bb$ e i moduli $|\ab|$, $|\bb|$.
	(b) Determinare l'angolo $\alpha$ formato dai due vettori.
	(c) Calcolare il prodotto vettoriale $\ab \times \bb$ e verificare le proprietà $\ab \times \bb = -\bb \times \ab$ e $\ab \cdot (\ab \times \bb) = 0$.
	(d) Verificare l'identità $|\ab \times \bb|^2 = |\ab|^2 |\bb|^2 - (\ab \cdot \bb)^2$.
	
	\risultato (a) $\ab \cdot \bb = 3$, $|\ab| = \sqrt{6}$, $|\bb| = \sqrt{14}$. (b) $\cos\alpha = 3/\sqrt{84}$, da cui $\alpha \simeq 70{,}9^\circ$. (c) $\ab \times \bb = (5\;\,-5\;\,5)^\top$, $|\ab \times \bb| = 5\sqrt{3}$. (d) Verifica numerica: $|\ab \times \bb|^2 = 75$, $|\ab|^2 |\bb|^2 - (\ab \cdot \bb)^2 = 84 - 9 = 75$.
	
\end{esercizio}

\begin{esercizio}[Verifica del doppio prodotto vettoriale]
	\label{ex:A03_doppio_prodotto}
	
	Si considerino i tre vettori
	\begin{equation*}
		\bm{a} = (1\;\,0\;\,2)^\top, \qquad
		\bm{b} = (0\;\,1\;\,1)^\top, \qquad
		\bm{c}= (2\;\,1\;\,0)^\top.
	\end{equation*}
	(a) Calcolare $\ab \times (\bb \times \cb)$ procedendo per sostituzione diretta.
	(b) Verificare l'identità del doppio prodotto vettoriale $\ab \times (\bb \times \cb) = (\ab \cdot \cb) \, \bb - (\ab \cdot \bb) \, \cb$ calcolando separatamente $(\ab \cdot \cb) \, \bb$ e $(\ab \cdot \bb) \, \cb$ e mostrando che la differenza coincide col risultato del punto (a).
	(c) Calcolare il prodotto misto $\ab \cdot (\bb \times \cb)$ e verificarne l'invarianza rispetto a permutazione ciclica.
	
	\risultato (a) $\bb \times \cb = (-1\;\,2\;\,-2)^\top$, $\ab \times (\bb \times \cb) = (-4\;\,0\;\,2)^\top$. (b) $\ab \cdot \cb = 2$, $\ab \cdot \bb = 2$, da cui $(\ab \cdot \cb) \bb - (\ab \cdot \bb) \cb = 2(0\;\,1\;\,1)^\top - 2(2\;\,1\;\,0)^\top = (-4\;\,0\;\,2)^\top$. (c) $\ab \cdot (\bb \times \cb) = -1 + 0 - 4 = -5$; per permutazione ciclica $\bb \cdot (\cb \times \ab)$ e $\cb \cdot (\ab \times \bb)$ producono lo stesso valore $-5$.
	
\end{esercizio}

\begin{esercizio}[Momento polare di una forza e legge di variazione]
	\label{ex:A03_momento_polare}
	
	Una forza $\bm{f} = (3\;\,-2\;\,1)^\top$ è applicata nel punto $P = (1,\,2,\,0)$. Si considerino due poli, $O = (0,\,0,\,0)$ e $O' = (2,\,1,\,1)$.
	
	(a) Calcolare il momento polare $\bm{\upmu}(O) = \overrightarrow{OP} \times \fb$ e il momento polare $\bm{\upmu}(O') = \overrightarrow{O'P} \times \fb$.
	
	(b) Verificare la legge di variazione del momento al variare del polo:
	\begin{equation*}
		\bm{\upmu}(O') = \bm{\upmu}(O) + \overrightarrow{O'O} \times \fb.
	\end{equation*}
	
	(c) Verificare che il prodotto scalare $\fb \cdot \bm{\upmu}(O)$ è invariante rispetto al polo (cioè vale anche $\fb \cdot \bm{\upmu}(O')$).
	
	\risultato (a) $\overrightarrow{OP} = (1\;\,2\;\,0)^\top$, $\bm{\upmu}(O) = (2\;\,-1\;\,-8)^\top$; $\overrightarrow{O'P} = (-1\;\,1\;\,-1)^\top$, $\bm{\upmu}(O') = (-1\;\,-2\;\,-1)^\top$. (b) $\overrightarrow{O'O} = (-2\;\,-1\;\,-1)^\top$, $\overrightarrow{O'O} \times \fb = (-3\;\,-1\;\,7)^\top$; $\bm{\upmu}(O) + \overrightarrow{O'O} \times \fb = (-1\;\,-2\;\,-1)^\top$. (c) $\fb \cdot \bm{\upmu}(O) = 6 + 2 - 8 = 0$, $\fb \cdot \bm{\upmu}(O') = -3 + 4 - 1 = 0$. Entrambi nulli (forza e raggio vettore sono complanari nel piano $z=0$, dunque $\bm{\upmu}$ è perpendicolare a tale piano --- caso particolare).
	
\end{esercizio}

\begin{esercizio}[Momento assiale rispetto a una retta]
	\label{ex:A03_momento_assiale}
	
	Una forza $\fb = (0\;\,4\;\,2)^\top$ è applicata nel punto $P = (3,\,0,\,0)$. Si consideri la retta $r$ passante per l'origine $O$ e parallela al versore $\bm{a}_r = \frac{1}{\sqrt{3}}(1\;\,1\;\,1)^\top$.
	
	(a) Calcolare il momento polare $\bm{\upmu}(O)$ rispetto a $O$.
	
	(b) Calcolare il momento assiale $\mu_r$ rispetto alla retta $r$ come componente di $\bm{\upmu}(O)$ lungo $\ab_r$:
	\begin{equation*}
		\mu_r = \ab_r \cdot \bm{\upmu}(O).
	\end{equation*}
	
	(c) Verificare che il momento assiale è invariante rispetto alla scelta del polo lungo $r$. Si calcoli a tal fine $\bm{\upmu}(O')$ con $O' = (1,\,1,\,1)$ (punto su $r$) e si mostri che $\ab_r \cdot \bm{\upmu}(O') = \mu_r$.
	
	\textit{Suggerimento.} Per il punto (c), usare la legge di variazione del momento e osservare che $\overrightarrow{O'O} = -\sqrt{3}\,\ab_r$ è parallelo ad $\ab_r$; il prodotto vettoriale $\overrightarrow{O'O} \times \fb$ è dunque perpendicolare ad $\ab_r$ e si annulla nel prodotto scalare con $\ab_r$.
	
	\risultato (a) $\overrightarrow{OP} = (3\;\,0\;\,0)^\top$, $\bm{\mu}(O) = (0\;\,-6\;\,12)^\top$. (b) $\mu_r = \frac{1}{\sqrt{3}}(0 - 6 + 12) = \frac{6}{\sqrt{3}} = 2\sqrt{3}$. (c) $\overrightarrow{O'P} = (2\;\,-1\;\,-1)^\top$, $\bm{\mu}(O') = (2\;\,-4\;\,8)^\top$; $\ab_r \cdot \bm{\upmu}(O') = \frac{1}{\sqrt{3}}(2-4+8) = 2\sqrt{3}$. Coincide con $\mu_r$.
	
\end{esercizio}

\begin{esercizio}[Natura polare o assiale di alcune grandezze fisiche]
	\label{ex:A03_polari_assiali}
	
	Per ciascuna delle seguenti grandezze, stabilire se si tratta di un vettore polare o di un vettore assiale, motivando la risposta sulla base del comportamento sotto riflessione.
	
	(a) La velocità $\vb$ di un punto materiale.
	
	(b) La velocità angolare $\bm{\upomega}$ di un corpo rigido in rotazione.
	
	(c) Il momento angolare $\mathbf{L} = \rb \times m\vb$ di un punto materiale, con $\rb$ vettore posizione e $\vb$ velocità.
	
	(d) La forza $\fb$ agente su un punto materiale.
	
	(e) Il momento di una forza $\bm{\upmu} = \rb \times \fb$.
	
	(f) Il prodotto vettoriale di due vettori assiali, $\bm{\upomega}_1 \times \bm{\upomega}_2$.
	
	\textit{Suggerimento.} Si ricordi la regola di composizione: il prodotto vettoriale di due polari è assiale, di due assiali è assiale, di un polare e un assiale è polare.
	
	\risultato (a) Polare: la velocità è la derivata temporale di un vettore posizione, anch'esso polare. (b) Assiale: la velocità angolare è definita dalla rotazione attraverso la regola della mano destra, e il suo segno dipende dall'orientazione della terna di riferimento. (c) Assiale: prodotto vettoriale di due polari ($\rb$ e $\vb$). (d) Polare. (e) Assiale: prodotto vettoriale di due polari. (f) Assiale: prodotto vettoriale di due assiali.
	
\end{esercizio}

\begin{esercizio}[Sistema di forze nel piano: momenti polari]
	\label{ex:A03_piano_momenti}
	
	Nel piano $(O; x, y)$, si considerino tre forze applicate:
	\begin{itemize}[noitemsep,topsep=2pt]
		\item $\bm{f}_1 = (3\;\,0)^\top$ applicata in $P_1 = (0,\,2)$;
		\item $\bm{f}_2 = (0\;\,-2)^\top$ applicata in $P_2 = (3,\,0)$;
		\item $\bm{f}_3 = (-1\;\,1)^\top$ applicata in $P_3 = (2,\,2)$.
	\end{itemize}
	
	(a) Calcolare il momento polare di ciascuna forza rispetto all'origine $O$, espresso come scalare con segno (positivo antiorario).
	
	(b) Calcolare il momento polare risultante del sistema rispetto a $O$.
	
	(c) Verificare la legge di variazione del momento, calcolando il momento risultante rispetto al polo $O' = (1,\,1)$ e confrontandolo col risultato ottenuto applicando la legge $\mu(O') = \mu(O) + \overrightarrow{O'O} \times \fb$, dove $\fb$ è il risultante delle forze.
	
	\textit{Suggerimento.} Nel piano, il momento di una forza $\bm{f} = (X\;\,Y)^\top$ applicata in $P = (x, y)$ rispetto all'origine è lo scalare $\mu = x \, Y - y \, X$.
	
	\risultato (a) $\mu_1 = -6$, $\mu_2 = -6$, $\mu_3 = 4$. (b) $\mu(O) = -8$. (c) Risultante $\bm{f} = (2\;\,-1)^\top$; $\overrightarrow{O'O} = (-1\;\,-1)^\top$; il prodotto vettoriale planare $\overrightarrow{O'O} \times \fb = (-1)(-1) - (-1)(2) = 1 + 2 = 3$; quindi $\mu(O') = -8 + 3 = -5$. Verifica diretta: $\mu(O') = (x_1 - 1)Y_1 - (y_1 - 1)X_1 + \ldots = -3 - 4 + 2 = -5$.
	
\end{esercizio}

\begin{esercizio}[Verifica del comportamento sotto riflessione]
	\label{ex:A03_riflessione}
	
	Si consideri la riflessione rispetto al piano $xy$, che inverte l'asse $z$. Sotto questa trasformazione, le componenti di un vettore polare $\bm{p} = (p_x\;\,p_y\;\,p_z)^\top$ si trasformano in $\\bm{p}' = (p_x\;\,p_y\;\,-p_z)^\top$, mentre le componenti di un vettore assiale $\bm{a} = (a_x\;\,a_y\;\,a_z)^\top$ si trasformano in $\bm{a}' = (-a_x\;\,-a_y\;\,a_z)^\top$.
	
	Si considerino i vettori:
	\begin{equation*}
		\bm{u} = (1\;\,2\;\,3)^\top, \qquad
		\bm{v} = (2\;\,-1\;\,1)^\top, \qquad
		\wb = \ub \times \vb.
	\end{equation*}
	
	(a) Calcolare $\wb$.
	
	(b) Trasformare $\ub$ e $\vb$ sotto la riflessione, ottenendo $\ub'$ e $\vb'$, e calcolare $\wb' = \ub' \times \vb'$.
	
	(c) Verificare che $\wb'$ si trasforma secondo la regola dei vettori \emph{assiali}, cioè $\wb' = (-w_x\;\,-w_y\;\,w_z)^\top$, confermando che il prodotto vettoriale di due polari è assiale.
	
	\risultato (a) $\bm{w} = (5\;\,5\;\,-5)^\top$. (b) $\bm{u}' = (1\;\,2\;\,-3)^\top$, $\bm{v}' = (2\;\,-1\;\,-1)^\top$, $\wb' = \ub' \times \vb' = (-5\;\,-5\;\,-5)^\top$. (c) La regola di trasformazione assiale prevede $(-w_x\;\,-w_y\;\,w_z)^\top = (-5\;\,-5\;\,-5)^\top$. Coincide con $\wb'$ calcolato direttamente: il prodotto vettoriale è effettivamente assiale.
	
\end{esercizio}


